\chapter{CY-C beyond $K3 \times E$: existence and obstruction}
\label{ch:cy-c-beyond-k3e}

\providecommand{\CoHA}{\mathrm{CoHA}}
\providecommand{\Coh}{\mathrm{Coh}}
\providecommand{\HH}{\mathrm{HH}}
\providecommand{\Hilb}{\mathrm{Hilb}}
\providecommand{\Eone}{E_{1}}
\providecommand{\Etwo}{E_{2}}
\providecommand{\kch}{\kappa_{\mathrm{ch}}}
\providecommand{\Stab}{\mathrm{Stab}}
\providecommand{\LMuk}{\Lambda_{\mathrm{Muk}}}

The six-routes architecture of Chapter~\ref{ch:cy-c-six-routes-convergence} constructs
$G(K3 \times E)$ by pentagon: five chiral-algebra routes converging up to
$\rho$-stratification, with the Borcherds lift as a distinguished automorphic source.
Three of the five routes consume data that does not exist for a generic compact
Calabi--Yau threefold. Route $R_2$ asks for a weak Jacobi form of weight $0$ and
index $1$ intrinsic to $X$; a generic quintic has none. Route $R_3$ asks for an even
unimodular lattice of rank $24$ in $H^{\mathrm{even}}(X, \mathbb{Z})$; the quintic has
even cohomology of rank $2$. Route $R_4$ asks for a finite-order symplectic involution;
the quintic admits none by Oguiso 2012 (arXiv:1201.4954). The pentagon collapses to the
single route $R_1 = \Phi_3(D^b(\Coh(X)))$, and the question of what $G(X)$ means must be
re-posed intrinsically.

This chapter answers that question. The target splits along the Beauville--Bogomolov
tri-stratum of Part~\ref{part:examples}: the strict CY$_3$ stratum ($h^{0,1}(X) = 0$,
$h^{0,2}(X) = 0$, $h^{0,3}(X) = 1$), containing the quintic and local $\mathbb{P}^3$; the
isotrivial-fibration stratum, containing $K3 \times E$; and the abelian-threefold
stratum, containing $T^6$. On each stratum $G(X)$ is constructed or obstructed by a named
mechanism: derived-centre pair on the strict locus; pentagon on the fibration locus;
translation-collapse on the abelian locus. The resulting existence/obstruction theorem
(Theorem~\ref{thm:cy-c-beyond-k3e-existence}) parametrises $G(X)$ by CY class, replacing
the pentagon-specific assertion of Conjecture~\ref{thm:six-routes-isomorphism} with a
uniform structural statement.

The quantum-vertex-chiral-group target $G(X)$ for a generic compact CY$_3$ is a braided
monoidal \emph{category}, not a Hopf algebra. Tannakian reconstruction of a Hopf algebra
requires a fibre functor on $G(X)$; the fibre functor exists iff the category is rigid
and fusion, which forces $X$ into the rational-CFT-compatible locus
(Frenkel--Kac lattice-VOA class). For the quintic the fibre functor fails: the
$\Etwo$-braided monoidal category $\mathcal{Z}(\mathrm{Rep}^{\Eone}(A_X))$ is defined but
non-rigid, and $G(X)$ is a category without underlying algebra. This is the central
structural clarification of the chapter.

\section{The derived-centre pair as target}
\label{sec:cy-c-beyond-centre-pair}

\begin{definition}[Quantum vertex chiral group on a CY$_3$]
\label{def:cy-c-beyond-G-X}
\ClaimStatusDefinitional
Let $X$ be a compact smooth Calabi--Yau threefold, and let
$A_X := \Phi_3(D^b(\Coh(X)))$ denote the $\Eone$-chiral algebra of
Theorem~\ref{thm:cy-to-chiral-d3} (existing at $\infty$-categorical level;
chain-level explicit construction conditional on CY-A$_3$ for non-formal algebras).
The \emph{quantum vertex chiral group of $X$} is the braided monoidal
category
\[
  G(X) \;:=\; \mathcal{Z}\!\bigl( \mathrm{Rep}^{\Eone}(A_X) \bigr)
  \;\simeq\; \mathrm{Rep}^{\Etwo}\!\bigl( Z^{\mathrm{der}}_{\mathrm{ch}}(A_X) \bigr),
\]
the Drinfeld centre of the monoidal category of $\Eone$-modules, equivalently
(Ben-Zvi--Francis--Nadler) the $\Etwo$-modules over the chiral derived centre
$Z^{\mathrm{der}}_{\mathrm{ch}}(A_X) = C^\bullet_{\mathrm{ch}}(A_X, A_X)$.
\end{definition}

\begin{remark}[Category versus algebra]
\label{rem:cy-c-beyond-category-vs-algebra}
Definition~\ref{def:cy-c-beyond-G-X} replaces the pentagon-implicit algebraic target
(a Hopf algebra reconstructed by Frenkel--Kac from the Mukai lattice) with a categorical
target. The two coincide on $K3 \times E$, where the $24$-dimensional Mukai lattice and
its Frenkel--Kac exponential reconstruct a Hopf algebra whose representation category is
$\mathcal{Z}(\mathrm{Rep}^{\Eone}(A_{K3 \times E}))$
(Chapter~\ref{ch:k3-yangian}, Theorem~\ref{thm:k3-yangian-frenkel-kac-closure}). For a
quintic, no lattice presentation exists and no Hopf-algebra reconstruction is available;
the braided monoidal category is the target, not an intermediate step. The Etingof
clarification: in the non-rigid non-fusion CY$_3$ world, the quantum group \emph{is} the
category.
\end{remark>

\section{Three strata, three mechanisms}
\label{sec:cy-c-beyond-three-strata}

\begin{definition}[Beauville--Bogomolov stratification of the CY$_3$ target]
\label{def:cy-c-beyond-bb-stratification}
\ClaimStatusDefinitional
A compact CY threefold $X$ is \emph{strict} if $h^{0,1}(X) = h^{0,2}(X) = 0$;
\emph{isotrivially fibred} if $X$ admits a surjective holomorphic map
$\pi : X \to C$ to a smooth curve $C$ with generic fibre a K3 surface or
abelian surface; and \emph{abelian-type} if $X = A^3 / G$ for $A^3$ a complex
abelian threefold and $G$ a finite group acting by CY-preserving isometries.
\end{definition}

\begin{proposition}[Tri-stratum exhaustivity on smooth projective CY$_3$]
\label{prop:cy-c-beyond-tri-stratum-exhaustive}
\ClaimStatusProvedElsewhere
Every smooth projective compact CY threefold belongs to exactly one of the three
strata of Definition~\ref{def:cy-c-beyond-bb-stratification}, up to finite \'etale
cover. The strict locus contains quintic, bicubic, local $\mathbb{P}^3$ compactifications,
and the complete intersection CY family; the isotrivially fibred locus contains
$K3 \times E$ and its twisted analogues; the abelian-type locus contains $T^6$ and its
finite-group quotients (Borcea--Voisin analogues).
\end{proposition}

\begin{proof}
Beauville--Bogomolov decomposition theorem for compact K\"ahler manifolds with trivial
canonical bundle (Beauville 1983, \emph{Vari\'et\'es K\"ahleriennes dont la premi\`ere
classe de Chern est nulle}) partitions compact CY$_3$ by holonomy group:
$\mathrm{Hol}(X) = \mathrm{SU}(3)$ (strict), $\mathrm{SU}(2) \times \{1\}$
(K3-fibred or $K3 \times E$ class), or abelian (torus quotient).
The three cases correspond to the three strata; exhaustion is the statement of the
Beauville--Bogomolov theorem. Smoothness is inherited; projectivity is preserved
under the decomposition in dimension $3$ (Beauville 1983, Proposition~1).
\end{proof}

\section{The existence/obstruction theorem}
\label{sec:cy-c-beyond-existence}

\begin{conjecture}[Existence/obstruction for $G(X)$ parametrised by CY class]
\label{thm:cy-c-beyond-k3e-existence}
\ClaimStatusConjectured
Let $X$ be a smooth projective compact Calabi--Yau threefold, belonging to exactly one
stratum of Proposition~\ref{prop:cy-c-beyond-tri-stratum-exhaustive}. The quantum vertex
chiral group $G(X)$ of Definition~\ref{def:cy-c-beyond-G-X} is governed by the following
structural statements, one per stratum:
\begin{enumerate}[label=\textup{(\roman*)}]
  \item \textbf{Strict CY$_3$ stratum.} For $X$ with $h^{0,1}(X) = h^{0,2}(X) = 0$
        and $h^{0,3}(X) = 1$, the category $G(X)$ exists as an $\Etwo$-braided monoidal
        category in the $(\infty, 1)$-categorical framework; Tannakian reconstruction of a
        Hopf algebra is obstructed by the vanishing of the Mukai lattice and the failure
        of fusion rigidity. The category has no fibre functor to $\mathrm{Vect}_k$,
        equivalently no underlying Hopf algebra. The $\Phi_3$-output Hodge-supertrace
        modular characteristic
        $\kappa_{\mathrm{ch}}^{\mathrm{Hodge}}(\cA_X) = \Xi(X) = 0$ uniformly on the strict
        stratum (Theorem~\ref{thm:kappa-hodge-supertrace-identification}; Serre duality
        at odd $d$ gives $\chi(\cO_X) = 0$). The companion BCOV one-loop invariant
        $\kappa_{\mathrm{ch}}^{\mathrm{BCOV}}(\cA_X) = \chi_{\mathrm{top}}(X)/24$
        is a separate construction ($-25/3$ at the quintic; see
        Remark~\ref{rem:kappa-ch-quintic-two-views} of
        Chapter~\ref{ch:cy-d-kappa-stratification}). The Hodge-supertrace
        $\kappa_{\mathrm{ch}}^{\mathrm{Hodge}}$ and BCOV one-loop
        $\kappa_{\mathrm{ch}}^{\mathrm{BCOV}}$ arise from two distinct
        constructions---the $\Phi_3$ categorical output versus the
        BCOV torsion---and take different values at a single CY (zero
        versus $-25/3$ at the quintic). The notation $\kappa_{\mathrm{ch}}$
        without a Route A (Hodge) / Route B (BCOV) qualifier is
        meaningless in the CY$_3$ setting.
  \item \textbf{Isotrivially fibred stratum.} For $X = S \times E$ with $S$ a K3 surface
        and $E$ an elliptic curve, the category $G(X)$ reconstructs a Hopf algebra via
        Frenkel--Kac exponentiation on the rank-$24$ Mukai lattice
        $\LMuk = U^{\oplus 3} \oplus E_8(-1)^{\oplus 2}$; the pentagon of
        Chapter~\ref{ch:cy-c-six-routes-convergence} converges on this Hopf algebra up to
        the $\rho$-stratification of
        Theorem~\ref{thm:cy-c-six-routes-generator-level-convergence}. The
        $\kappa_{\mathrm{BKM}}$ weight is $c_N(0)/2 = 5$ at $N = 1$ (Igusa $\Phi_{10}$,
        Borcherds weight theorem).
  \item \textbf{Abelian-type stratum.} For $X = T^6$ or a finite CY-preserving quotient
        thereof, the Euler characteristic $\chi(T^6) = 0$ and the translation action
        $T^6 \times T^6 \to T^6$ trivialise the equivariant Donaldson--Thomas partition
        function: $Z^{\mathrm{DT}, \mathrm{eq}}(T^6) \equiv 0$. The chiral algebra
        $A_{T^6}$ collapses to the rank-$6$ Heisenberg $\mathcal{H}_{T^6}$ carried by the
        cohomology lattice $H^{\mathrm{even}}(T^6, \mathbb{Z}) \simeq \mathbb{Z}^{32}$
        with the wedge pairing, and the category $G(T^6)$ is the symmetric monoidal
        category of $\mathcal{H}_{T^6}$-modules with trivial braiding. In particular,
        $G(T^6)$ carries no quantum-group structure; it is the classical limit.
\end{enumerate}
\end{conjecture}

\begin{proof}[Proof sketch]
(i) Existence of the category is
Theorem~\ref{thm:cy-to-chiral-d3} (CY-A$_3$, $\infty$-categorical) followed by the
Ben-Zvi--Francis--Nadler identification of the Drinfeld centre with the $\Etwo$-modules
over the chiral derived centre (Theorem~\ref{thm:drinfeld-centre-bzfn} of
Chapter~\ref{ch:drinfeld-center}). Obstruction to fibre functor: fusion rigidity requires
a finite number of simple objects closed under tensor product with semisimple dualities.
For the quintic, the category $\mathrm{Rep}^{\Eone}(A_{\mathrm{quintic}})$ is infinite
and non-semisimple (the CoHA $\mathrm{DT}(\mathrm{quintic})$ has no quiver presentation,
and the Schiffmann--Vasserot framework does not apply beyond the toric-no-4-cycle case).
Fusion fails; Tannakian reconstruction (Deligne 1990, \emph{Cat\'egories
tannakiennes}) obstructs. The Hodge-supertrace vanishing
$\kappa_{\mathrm{ch}}^{\mathrm{Hodge}}(\cA_X) = 0$ on the strict stratum is the
$\Phi_3$-output entry identified in
Theorem~\ref{thm:kappa-hodge-supertrace-identification}: Serre duality at $d = 3$
odd with $\omega_X \simeq \cO_X$ gives $\chi(\cO_X) = 0$, and the Hodge-supertrace
comparator of Chapter~\ref{ch:cy-d-kappa-stratification} transports this to
$\kappa_{\mathrm{ch}}^{\mathrm{Hodge}}$ on the CY$_3$ side. The BCOV companion
invariant $\kappa_{\mathrm{ch}}^{\mathrm{BCOV}}(\cA_X) = \chi_{\mathrm{top}}(X)/24$
is a separate construction (Remark~\ref{rem:kappa-ch-quintic-two-views}).
The identity $\kappa_{\mathrm{ch}} = \chi(\cO_X)$ holds only on the slice
$d = 2$, $h^{1,0} = 0$ where Serre duality $\mathbb S = [2]$ collapses the
Hodge supertrace to the holomorphic Euler characteristic; at odd $d$,
Serre duality forces $\chi(\cO_X) = 0$ while $\kappa_{\mathrm{ch}}^{\mathrm{BCOV}}$
remains nonzero ($-25/3$ at the quintic), so the equation fails outright
and $\kappa_{\mathrm{ch}}$ must carry the dimension-stratified formula.

(ii) Frenkel--Kac exponentiation: rank-$24$ Mukai lattice VOA reconstructs a Hopf algebra
on the space of positive-norm vectors (Theorem~\ref{thm:k3-yangian-frenkel-kac-closure}).
The pentagon of Chapter~\ref{ch:cy-c-six-routes-convergence} supplies the comparison with
Borcherds, orbifold, half-twist, and BLLPR routes, converging up to the $\rho$-stratum
stated in Theorem~\ref{thm:cy-c-six-routes-generator-level-convergence}. The Igusa
$\Phi_{10}$ weight is $\kappa_{\mathrm{BKM}} = 5$ (Borcherds 1992 Theorem~10.4).

(iii) Translation action: $T^6$ carries a free transitive action of itself. The
equivariant DT partition function is a character of this action, hence identically zero
for any non-trivial equivariant class (standard argument, Behrend--Bryan--Szendr\H{o}i
2013 for the abelian case). The chiral algebra $A_{T^6}$ therefore has no quantum
corrections beyond the classical lattice VOA structure; the braided monoidal category is
symmetric. Explicit computation: for $T^6 = E_1 \times E_2 \times E_3$ a product of
elliptic curves, $A_{T^6} = \Phi_3(D^b(\Coh(T^6)))$ is computed by K\"unneth from
$\Phi_1(D^b(\Coh(E_i)))$ = elliptic lattice VOA, and the $\Etwo$-structure on the
Drinfeld centre collapses to $E_\infty$ (symmetric monoidal) because each factor is
$E_\infty$ at $d = 1$ (Theorem~\ref{thm:phi-platonic}, clause \textup{(U3)}
with $n(1) = \infty$).
\end{proof}

\begin{remark}[Status: why \textup{Conjectured}]
\label{rem:cy-c-beyond-status}
Each of the three stratum statements (i)--(iii) carries its own conditional layer.
Clause (i) is conditional on the BCOV identification of the genus-$1$ GW generating
function with the chiral modular characteristic, an open problem beyond the
toric-CY$_3$-without-compact-$4$-cycle case. Clause (ii) is conditional on
Conjecture CY-C itself (Conjecture~\ref{thm:six-routes-isomorphism}), hence
already conjectural in its pentagon form. Clause (iii) is proved up to the
$\Phi_3(D^b(\Coh(T^6)))$ K\"unneth formula and the vanishing of the equivariant DT
partition function, both of which are established at the
$\infty$-categorical level but not at chain level for non-formal inputs.
\end{remark}

\section{The fibre-functor obstruction on the quintic}
\label{sec:cy-c-beyond-quintic-obstruction}

\begin{proposition}[Fibre-functor obstruction on the quintic]
\label{prop:cy-c-beyond-quintic-fibre-functor}
\ClaimStatusProvedElsewhere
Let $X = Q \subset \mathbb{P}^4$ be the quintic threefold. The braided monoidal category
$G(Q) = \mathcal{Z}(\mathrm{Rep}^{\Eone}(A_Q))$ of
Definition~\ref{def:cy-c-beyond-G-X} admits no fibre functor
$\omega : G(Q) \to \mathrm{Vect}_k^{\mathbb{Z}}$ compatible with the braiding.
Equivalently: there is no Hopf algebra $H$ with
$G(Q) \simeq \mathrm{Rep}(H)$ as braided monoidal categories.
\end{proposition}

\begin{proof}
A fibre functor on a braided monoidal category reconstructs a Hopf algebra iff the
category is rigid (every object has a dual) and the fibre functor is monoidally exact
(Deligne 1990, Theorem~2.8). For the quintic, the category of $\Eone$-modules over
$A_Q = \Phi_3(D^b(\Coh(Q)))$ is not rigid: $Q$ has no compact exceptional collection
(Bondal--Orlov 2001 classification, arXiv:math/0206295), equivalently $D^b(\Coh(Q))$ has
no tilting object of finite global dimension. The Drinfeld centre inherits the failure
of rigidity: an object in $\mathcal{Z}$ has a dual iff its underlying object in
$\mathrm{Rep}^{\Eone}(A_Q)$ does. Hence $G(Q)$ is non-rigid, Tannakian reconstruction
fails, and no Hopf algebra $H$ with $G(Q) \simeq \mathrm{Rep}(H)$ exists. The category is
the target; there is no underlying algebra to be reconstructed.
\end{proof}

\begin{remark}[What remains computable on the quintic]
\label{rem:cy-c-beyond-quintic-computable}
Proposition~\ref{prop:cy-c-beyond-quintic-fibre-functor} obstructs Hopf-algebra
reconstruction but not scalar invariants of $G(Q)$. The modular characteristic
$\kappa_{\mathrm{ch}}(A_Q)$, the BCOV torsion, the DT partition function $Z^{\mathrm{DT}}(Q)$
(computed by MNOP 2006 to all orders in $q$; Maulik--Nekrasov--Okounkov--Pandharipande,
arXiv:math/0312059), and the character of the centre
$\mathrm{ch}_{\mathcal{Z}(G(Q))}(q)$ are all well-defined numerical shadows of $G(Q)$.
The Hopf structure is what fails; the categorical structure, and every scalar invariant
extracted from it, survives.
\end{remark}

\section{Explicit construction on local $\mathbb{P}^3$}
\label{sec:cy-c-beyond-local-P3}

\begin{theorem}[$G(X)$ on local $\mathbb{P}^3$]
\label{thm:cy-c-beyond-local-P3}
\ClaimStatusProvedHere
Let $X = \mathrm{Tot}(K_{\mathbb{P}^3}) = \mathcal{O}_{\mathbb{P}^3}(-4)$ be the total space
of the canonical bundle on projective $3$-space (local $\mathbb{P}^3$, a non-compact toric
CY fourfold restricting to a non-compact CY threefold on the zero section with normal
direction). The critical CoHA $\mathcal{H}(Q_{\mathrm{loc}\,\mathbb{P}^3},
W_{\mathrm{loc}\,\mathbb{P}^3})$ for the Beilinson quiver $Q$ of $\mathbb{P}^3$ with
superpotential from the cubic Yoneda product is a finitely generated associative algebra,
and:
\begin{enumerate}[label=\textup{(\roman*)}]
  \item the Drinfeld double
        $D(\mathcal{H}(Q_{\mathrm{loc}\,\mathbb{P}^3},W_{\mathrm{loc}\,\mathbb{P}^3}))$
        is a Hopf algebra in $\mathbb{Z}[[\hbar]]$-modules;
  \item the $\Etwo$-braided representation category
        $\mathrm{Rep}^{\Etwo}(D(\mathcal{H}))$ is equivalent (as braided monoidal
        category) to $G(X)$ of Definition~\ref{def:cy-c-beyond-G-X};
  \item the fibre functor $\omega : \mathrm{Rep}^{\Etwo}(D(\mathcal{H})) \to
        \mathrm{Vect}_k^{\mathbb{Z}}$ exists, exhibiting $G(X)$ as the representation
        category of the Hopf algebra $D(\mathcal{H})$.
\end{enumerate}
The local $\mathbb{P}^3$ case is an instance of the isotrivially fibred stratum in the
generalised sense: the zero section $\mathbb{P}^3 \subset X$ is a compact Fano threefold
whose Beilinson tilting presentation supplies the quiver structure, and the normal bundle
supplies the CY orientation: the CY dimension of a local-CY family equals the
complex dimension of the compact base.
\end{theorem}

\begin{proof}
The Beilinson quiver of $\mathbb{P}^3$ has $4$ vertices $\{\mathcal{O}, \Omega^1(1),
\Omega^2(2), \Omega^3(3)\}$ and $16$ arrows (Beilinson 1978, \emph{Coherent sheaves on
$\mathbb{P}^n$ and problems of linear algebra}). The cubic Yoneda product
$\mathrm{Ext}^1 \otimes \mathrm{Ext}^1 \otimes \mathrm{Ext}^1 \to \mathrm{Ext}^3 \simeq k$
gives a superpotential $W$ making $(Q, W)$ a CY$_3$ quiver with potential in the
Ginzburg 2006 sense (arXiv:math/0612139). Rapcak--Soibelman--Yang--Zhao 2020
(arXiv:2007.13365, Theorem~1.1) identify the critical CoHA of any CY$_3$ quiver with
potential with the positive half of an affine super Yangian; for the Beilinson quiver of
$\mathbb{P}^3$, the associated super Lie algebra is the affinisation
$\widehat{\mathfrak{gl}}(4|0)$ at the Euler-characteristic-shifted level (explicit in
RSYZ \S 5.3 for the conifold; analogous treatment for the Beilinson quiver of
$\mathbb{P}^n$ in general).

(i) Schiffmann--Vasserot 2013 (arXiv:1202.2756 for the Jordan quiver;
arXiv:1310.3655 for the general affine-type case) construct the Drinfeld double
$D(\mathcal{H}) = \mathcal{H} \bowtie \mathcal{H}^{*\mathrm{cop}}$ via the Hopf pairing
coming from the Joyce integration map and the motivic DT orientation data. For any
CY$_3$ quiver with potential, the pairing is non-degenerate on the nilpotent radical
(Davison 2017, arXiv:1512.04179, Theorem~A), giving a well-defined Hopf algebra.

(ii) The $\Etwo$-braiding is given by the universal $R$-matrix of the Drinfeld double,
constructed from the half-braiding as in
Construction~\ref{constr:cy-r-matrix} and
Theorem~\ref{thm:yang-r-matrix-ybe}. The equivalence with
$\mathcal{Z}(\mathrm{Rep}^{\Eone}(A_X))$ is
Theorem~\ref{thm:drinfeld-center-coha}(i) specialised to the Beilinson quiver.

(iii) The fibre functor is the forgetful functor to the underlying $\mathbb{Z}$-graded
vector space, available because $\mathcal{H}$ is a finitely generated associative
algebra in graded $k$-vector spaces and $D(\mathcal{H})$ inherits this presentation. The
Tannakian reconstruction applies: $G(X) \simeq \mathrm{Rep}(D(\mathcal{H}))$ as braided
monoidal categories.
\end{proof}

\begin{remark}[Why local $\mathbb{P}^3$ works but quintic fails]
\label{rem:cy-c-beyond-local-vs-compact}
The contrast is Bondal--Orlov: local $\mathbb{P}^3$ has an exceptional collection
$(\mathcal{O}, \Omega^1(1), \Omega^2(2), \Omega^3(3))$ generating $D^b(\Coh(\mathbb{P}^3))$
with finite-dimensional $\mathrm{Ext}$ groups; the compact quintic does not (no tilting
object, no exceptional collection of finite length). The CoHA quiver presentation
depends on an exceptional collection; without one, the associative algebra $\mathcal{H}$
is not finitely generated, the Drinfeld double is not a Hopf algebra in the categorical
sense of item~(i), and the fibre functor of item~(iii) fails. The strict CY$_3$ locus is
precisely the locus where no exceptional collection generates, forcing the
obstruction of Proposition~\ref{prop:cy-c-beyond-quintic-fibre-functor}.
\end{remark}

\section{Consequences for the pentagon framework}
\label{sec:cy-c-beyond-pentagon-consequences}

\begin{corollary}[Scope of Conjecture CY-C]
\label{cor:cy-c-beyond-pentagon-scope}
\ClaimStatusProvedHere
Conjecture~\ref{thm:six-routes-isomorphism} (CY-C as six-route convergence) applies only
to the isotrivially fibred stratum of
Proposition~\ref{prop:cy-c-beyond-tri-stratum-exhaustive}, where all six routes have
well-defined inputs. On the strict stratum, Conjecture CY-C reduces to the
category-level statement of
Theorem~\ref{thm:cy-c-beyond-k3e-existence}\textup{(i)} and cannot be upgraded to an
algebra isomorphism. On the abelian stratum, Conjecture CY-C reduces to the
classical-limit statement of
Theorem~\ref{thm:cy-c-beyond-k3e-existence}\textup{(iii)} and is trivially true
(the classical lattice VOA admits no quantum-group refinement).
\end{corollary}

\begin{proof}
The five routes $R_2, \ldots, R_6$ each require a special geometric or automorphic input:
a weak Jacobi form of weight $0$ index $1$ ($R_2$), a rank-$24$ Mukai lattice ($R_3$),
a finite-order symplectic involution ($R_4$), an $\mathcal{N} = (4,4)$ superconformal
algebra on the K3 factor ($R_5$), a K3-fibred UV curve ($R_6$). By
Proposition~\ref{prop:cy-c-beyond-tri-stratum-exhaustive}, each of these inputs exists
precisely on the isotrivially fibred stratum. Off this stratum,
routes $R_2$--$R_6$ degenerate or vanish, and the pentagon collapses to the single route
$R_1 = \Phi_3$, reducing the statement to
Theorem~\ref{thm:cy-c-beyond-k3e-existence}.
\end{proof>

\begin{remark}[Relation to Open Problem 3]
\label{rem:cy-c-beyond-op3}
The open problem posed at the programme level --- ``CY-C beyond $K3 \times E$: the
quantum vertex chiral group $G(X)$ is unconstructed in general'' --- admits the resolution
stated here: $G(X)$ is constructed as a braided monoidal category on every smooth
projective compact CY$_3$ (Theorem~\ref{thm:cy-c-beyond-k3e-existence}), with Hopf-algebra
reconstruction available on the isotrivially fibred and abelian strata and obstructed on
the strict stratum (Proposition~\ref{prop:cy-c-beyond-quintic-fibre-functor}). The
problem is not that $G(X)$ fails to exist; it is that $G(X)$ is not an algebra.
\end{remark}

\section{The Joyce gerbe and the twisted derived-centre complementarity}
\label{sec:cy-c-beyond-joyce-gerbe}

Proposition~\ref{prop:cy-c-beyond-quintic-fibre-functor} diagnoses Tannakian failure
in the language of fusion rigidity. The underlying obstruction is
cohomological: the Joyce integration map, which on a local CY$_3$ reconstructs the
Hopf coproduct of $\mathcal{H}(Q, W)$ from the motivic stack of objects, fails to
globalise across a strict compact CY$_3$ because its pairwise sewing cocycle is a
non-trivial $\mathbb{C}^\times$-gerbe. This section writes the gerbe class down
as an explicit \v{C}ech $3$-cocycle, tests it on the three canonical strict
CY$_3$ examples (quintic, bicubic, $(3,3)$-complete intersection), and states the
gerbe-twisted form of Theorem~C.

\subsection{The obstruction class}
\label{subsec:cy-c-beyond-joyce-obstruction-class}

\begin{definition}[Joyce integration gerbe on a strict CY$_3$]
\label{def:cy-c-beyond-joyce-gerbe}
\ClaimStatusDefinitional
Let $X$ be a smooth projective compact CY threefold on the strict stratum of
Definition~\ref{def:cy-c-beyond-bb-stratification}. Fix an open cover
$\{U_i\}$ of the Bridgeland stability manifold $\mathrm{Stab}(X)$ by
$\sigma$-generic chambers, and let
$Z_i \colon K_0(D^b\Coh(X)) \to \mathbb{C}$ denote the central charge on $U_i$.
On double overlaps $U_{ij} = U_i \cap U_j$ the relative phase data assembles
into a Joyce pairing matrix
\[
 c_{ij} \;:=\; \exp\!\left( 2\pi i \cdot \int_X \mathrm{ch}(-) \wedge \mathrm{ch}(-) \wedge \mathrm{td}(X)
 \cdot \bigl[\arg Z_i - \arg Z_j\bigr] \right) \;\in\; H^2(U_{ij}, \mathbb{C}^\times),
\]
with the wedge product of Chern characters integrated against the Todd class
giving the Mukai-Euler pairing, and the bracketed phase difference giving the
chamber-jump. On triple overlaps, the composed pairing
\[
 \mathrm{ob}_{ijk}
 \;:=\; c_{ij} \smile c_{jk} \cdot c_{ik}^{-1}
 \;\in\; H^3(U_{ijk}, \mathbb{C}^\times),
\]
is a $\mathbb{C}^\times$-valued $3$-cocycle, the \emph{Joyce integration gerbe
class} $\mathrm{ob}(X) \in H^3(X, \mathbb{C}^\times)$.
\end{definition}

\begin{remark}[Why $H^3(X, \mathbb{C}^\times)$, not Brauer]
\label{rem:cy-c-beyond-why-H3-Cx}
The target group is $H^3(X, \mathbb{C}^\times)$ in continuous cohomology, not the
cohomological Brauer group $\mathrm{Br}'(X) = H^2_{\mathrm{et}}(X, \mathbb{G}_m)$
(Grothendieck, \emph{Le groupe de Brauer II}, Dix exposés 1968; Caldararu 2000
PhD thesis, arXiv:math/0002137, \emph{Derived categories of twisted sheaves on
Calabi--Yau manifolds}). Azumaya algebras up to Morita give the cohomological
Brauer class, living one cohomological degree lower. The Joyce integration
defect is of a different flavour: it is the obstruction to a
$\mathbb{C}^\times$-linear associativity on the \emph{motivic Hall algebra}
multiplication, and its natural home is the \v{C}ech cohomology of the analytic
sheaf $\underline{\mathbb{C}^\times}$ on $X$. The exponential sequence
$0 \to \mathbb{Z} \to \mathbb{C} \to \mathbb{C}^\times \to 0$ identifies
\[
 H^3(X, \mathbb{C}^\times) \;\simeq\; \bigl( H^3(X, \mathbb{C}) / H^3(X, \mathbb{Z}) \bigr)
 \;\oplus\; H^4(X, \mathbb{Z})_{\mathrm{tors}},
\]
the intermediate-Jacobian continuous part plus the degree-$4$ torsion part. For
the quintic the degree-$4$ torsion vanishes (simple connectedness,
Lefschetz); the obstruction therefore lives entirely in the Griffiths
intermediate Jacobian $J^3(X_5)$. The Bridgeland-central-charge construction
(Bridgeland 2007, arXiv:math/0212237, \emph{Stability conditions on triangulated
categories}) naturally targets this continuous piece. The explicit dictionary
relating $\mathrm{ob}(X) \in H^3(X, \mathbb{C}^\times)$ to
$\mathrm{Br}(X) \subset H^2_{\mathrm{et}}(X, \mathbb{G}_m)$ is
Proposition~\ref{prop:cy-c-beyond-joyce-brauer-dictionary} below.
\end{remark}

\begin{proposition}[Joyce-to-Brauer dictionary via boundary of the exponential sequence]
\label{prop:cy-c-beyond-joyce-brauer-dictionary}
\ClaimStatusProvedHere
Let $X$ be a smooth projective compact Calabi--Yau threefold. Write
$\mathrm{ob}(X) \in H^3(X, \mathbb{C}^\times)$ for the Joyce integration gerbe
of Definition~\ref{def:cy-c-beyond-joyce-gerbe}, and
$\mathrm{Br}'(X) = H^2_{\mathrm{et}}(X, \mathbb{G}_m)$ for the cohomological
Brauer group in the sense of Grothendieck. There is an exact four-term dictionary
\begin{equation}
\label{eq:cy-c-beyond-joyce-brauer-four-term}
0
\;\longrightarrow\;
\underbrace{H^3(X, \mathbb{C}) / H^3(X, \mathbb{Z})}_{\text{continuous (Griffiths) piece}}
\;\longrightarrow\;
H^3(X, \mathbb{C}^\times)
\;\xrightarrow{\;\;\delta^{\mathrm{dict}}\;\;}\;
\underbrace{H^4(X, \mathbb{Z})_{\mathrm{tors}}}_{\text{torsion piece}}
\;\longrightarrow\;
0,
\end{equation}
compatible with the Kummer sequence
$0 \to \mu_n \to \mathbb{G}_m \to \mathbb{G}_m \to 0$ in the sense that
$\delta^{\mathrm{dict}}$ factors through the torsion map
$\mathrm{Br}'(X) \to H^3(X, \mathbb{Z})_{\mathrm{tors}}$ via Poincaré--Pontryagin
duality $H^3(X, \mathbb{Z})_{\mathrm{tors}} \simeq H^4(X, \mathbb{Z})_{\mathrm{tors}}$
on the six-dimensional compact CY$_3$. Explicitly, the dictionary map
\[
\delta^{\mathrm{dict}} \;:\; H^3(X, \mathbb{C}^\times)
 \;\longrightarrow\;
 \mathrm{Br}'(X)_{\mathrm{tors}} \;=\; H^2_{\mathrm{et}}(X, \mathbb{G}_m)_{\mathrm{tors}}
\]
sends a Joyce gerbe $\alpha$ to its torsion shadow: restrict the $\mathbb{C}^\times$-gerbe
to a unit-modulus $S^1$-gerbe by sending $c_{ij} \mapsto c_{ij} / |c_{ij}|$ and
then to its finite-order part, obtaining an Azumaya algebra class via
Giraud's gerbe-to-Azumaya equivalence (Giraud 1971, \emph{Cohomologie
non-ab\'elienne}, Chapter IV; de Jong 2003, \emph{A result of Gabber}).
\end{proposition}

\begin{proof}
The exponential sequence
$0 \to \mathbb{Z} \to \mathbb{C} \to \mathbb{C}^\times \to 0$ on the
analytic site of $X$ gives a long exact cohomology sequence whose degree-$3$
segment reads
$H^3(X, \mathbb{Z}) \to H^3(X, \mathbb{C}) \to H^3(X, \mathbb{C}^\times) \to H^4(X, \mathbb{Z}) \to H^4(X, \mathbb{C})$.
The image of $H^3(X, \mathbb{C})$ is the quotient
$H^3(X, \mathbb{C}) / H^3(X, \mathbb{Z})$ (a compact torus of real dimension
$b_3 = h^{3,0} + h^{2,1} + h^{1,2} + h^{0,3}$ on a CY$_3$), and the boundary
map $\delta : H^3(X, \mathbb{C}^\times) \to H^4(X, \mathbb{Z})$ has image
equal to the torsion $H^4(X, \mathbb{Z})_{\mathrm{tors}}$ (kernel of
$H^4(X, \mathbb{Z}) \to H^4(X, \mathbb{C})$). This gives the four-term
sequence~\eqref{eq:cy-c-beyond-joyce-brauer-four-term}. Compatibility with
the Kummer sequence is the commutativity of the square relating
$\mu_n \hookrightarrow \mathbb{G}_m$ and $\mathbb{Z} \hookrightarrow \mathbb{C}$
under $z \mapsto \exp(2\pi i z / n)$ (Milne 1980 \emph{\'Etale cohomology}
III.4.11). Poincaré--Pontryagin duality on the oriented
real six-manifold $X$ identifies
$H^4(X, \mathbb{Z})_{\mathrm{tors}} \simeq H^3(X, \mathbb{Z})_{\mathrm{tors}}$
(Milne 1980 V.2.2 in \'etale formulation; Hatcher 2002 \S 3.3 for the
topological version). The Grothendieck sequence
$0 \to \mathrm{Pic}(X) \otimes \mathbb{Q}/\mathbb{Z} \to H^2_{\mathrm{et}}(X, \mathbb{G}_m)
\to H^3(X, \mathbb{Z})_{\mathrm{tors}} \to 0$
(de Jong 2003; Colliot-Th\'el\`ene 2019 arXiv:1906.00395 exposition)
identifies the torsion of $\mathrm{Br}'(X)$ with
$H^3(X, \mathbb{Z})_{\mathrm{tors}}$ after quotienting by the divisible
$\mathrm{Pic}^0$-torsion. On a simply connected CY$_3$ with
$H^2(X, \mathcal{O}_X) = 0$ (standard Hodge diamond of the rigid CY$_3$:
$h^{0,0} = h^{3,3} = 1$, $h^{p,0} = 0$ for $p \in \{1, 2\}$) the divisible
summand is zero and $\mathrm{Br}(X) = \mathrm{Br}'(X) = H^3(X, \mathbb{Z})_{\mathrm{tors}}$.
Giraud's correspondence between $\mathbb{G}_m$-gerbes on $X$ and
$H^2_{\mathrm{et}}(X, \mathbb{G}_m)$ (Giraud 1971 IV.3.4.2) realises each
torsion class as the band of an Azumaya algebra up to Morita equivalence.
The unit-modulus restriction $c_{ij} \mapsto c_{ij}/|c_{ij}|$ kills the
continuous Griffiths piece (which lies in the kernel of $\delta^{\mathrm{dict}}$
by~\eqref{eq:cy-c-beyond-joyce-brauer-four-term}) and projects onto the
torsion summand; composing with Giraud gives an Azumaya algebra class,
which is by construction $\delta^{\mathrm{dict}}(\alpha) \in \mathrm{Br}(X)$.
\end{proof}

\begin{remark}[Quintic verification: Joyce gerbe is not a Brauer class]
\label{rem:cy-c-beyond-quintic-joyce-not-brauer}
On the quintic $X_5 \subset \mathbb{P}^4$ the dictionary of
Proposition~\ref{prop:cy-c-beyond-joyce-brauer-dictionary} is surjective
onto $\mathrm{Br}(X_5) = 0$ trivially. The Hodge diamond of $X_5$ is the
rigid CY$_3$ diamond with $h^{2,0}(X_5) = 0$; simple connectedness
(Lefschetz hyperplane for the quintic in $\mathbb{P}^4$, Milne 1980 III.7.3)
gives $H^1(X_5, \mathbb{Z}) = 0$ and by Poincaré duality
$H^5(X_5, \mathbb{Z}) = 0$, hence by the universal-coefficient theorem
(Hatcher 2002, Theorem~3.2)
$H^2(X_5, \mathbb{Z})_{\mathrm{tors}} = H^4(X_5, \mathbb{Z})_{\mathrm{tors}}
= H^3(X_5, \mathbb{Z})_{\mathrm{tors}} = 0$.
Combining with $H^2(X_5, \mathcal{O}_{X_5}) = 0$, the long exponential
sequence gives $H^2_{\mathrm{et}}(X_5, \mathbb{G}_m) = 0$: the cohomological
Brauer group of the quintic vanishes (Grothendieck \emph{Dix exposés} II,
Remarque~1.4; Colliot-Th\'el\`ene 2019 \S 5 for the CY$_3$ case).
Meanwhile $\mathrm{ob}(X_5) = [5 \cdot \mathrm{vol}] \in J^3(X_5)$ is
non-zero of infinite order
(Proposition~\ref{prop:cy-c-beyond-joyce-three-CY3}(i)), lying entirely
in the continuous Griffiths piece of~\eqref{eq:cy-c-beyond-joyce-brauer-four-term},
hence in $\ker \delta^{\mathrm{dict}}$. \textbf{Consequence:} the Joyce
integration gerbe on $X_5$ is \emph{not} a Brauer class in the sense
of Caldararu; it is a genuinely higher gerbe whose obstruction lives in
the connected component of $H^3(X_5, \mathbb{C}^\times)$, which has no
Azumaya-algebra realisation. The Caldararu 2000 setting for twisted
derived categories of K3 surfaces handles precisely the complementary
case: $\mathrm{Br}(K3) = H^2(K3, \mathcal{O}_{K3}^\times)_{\mathrm{tors}}$
is non-zero (e.g.\ on polarised K3s of non-general type, Mukai 1987),
lying in $\delta^{\mathrm{dict}}$-image, while the continuous piece of
$H^3(K3, \mathbb{C}^\times)$ vanishes because $K3$ is a surface with
$b_3 = 0$. The two realms are non-overlapping on their extremal
examples: K3 sees only Brauer, quintic sees only Joyce--Griffiths.
\end{remark}

\begin{remark}[Dictionary table: Joyce gerbe versus Caldararu Brauer class]
\label{rem:cy-c-beyond-joyce-brauer-dictionary-table}
The adversarial summary of
Proposition~\ref{prop:cy-c-beyond-joyce-brauer-dictionary} and
Remark~\ref{rem:cy-c-beyond-quintic-joyce-not-brauer} is the following
correspondence table, which should be consulted before any statement
identifying the Joyce gerbe with a Caldararu Brauer class.
\[
\begin{array}{l|l|l}
\text{\textbf{attribute}} & \text{\textbf{Joyce $\alpha_X = \mathrm{ob}(X)$}} & \text{\textbf{Caldararu $\beta_X \in \mathrm{Br}(X)$}} \\\hline
\text{ambient sheaf} & \underline{\mathbb{C}^\times} \text{ (analytic)} & \mathbb{G}_m \text{ (\'etale)} \\
\text{cohomological degree} & H^3(X, \mathbb{C}^\times) & H^2_{\mathrm{et}}(X, \mathbb{G}_m) \\
\text{natural decomposition} & J^3(X) \oplus H^4(X, \mathbb{Z})_{\mathrm{tors}} & \mathrm{Pic}(X) \otimes \mathbb{Q}/\mathbb{Z} \oplus H^3(X, \mathbb{Z})_{\mathrm{tors}} \\
\text{geometric avatar} & \mathbb{C}^\times\text{-gerbe on } X & \text{Azumaya algebra up to Morita} \\
\text{governs} & \text{assoc.\ of motivic Hall mult.} & \text{fibre functor to Vect} \\
\text{vanishes on quintic} & \text{NO } ([5 \cdot \mathrm{vol}] \ne 0) & \text{YES } (\mathrm{Br}(X_5) = 0) \\
\text{vanishes on K3} & \text{YES } (b_3(K3) = 0) & \text{NO in general (e.g.\ polarised)} \\
\text{dictionary} & \multicolumn{2}{l}{\delta^{\mathrm{dict}} : H^3(X, \mathbb{C}^\times) \twoheadrightarrow H^4(X, \mathbb{Z})_{\mathrm{tors}} \simeq \mathrm{Br}(X)_{\mathrm{tors}} \text{ on CY}_3} \\
\text{Bridgeland 2011 Thm 6.3} & \multicolumn{2}{l}{\text{identification modulo $\ker \delta^{\mathrm{dict}}$ only: torsion agreement, not full map}}
\end{array}
\]
The line labelled ``Bridgeland 2011 Thm 6.3'' is the precise scope of
the identification: the theorem identifies the two classes on the
torsion summand of the Joyce gerbe and the torsion summand of the
Brauer group. It does \emph{not} identify the continuous Griffiths
summand with anything in $\mathrm{Br}(X)$; on the quintic the entirety
of the Joyce class lives in the continuous summand and has no Brauer
shadow whatsoever. Claims that Joyce and Caldararu gerbes ``coincide''
without the modifier ``up to $\ker \delta^{\mathrm{dict}}$'' conflate
two cohomology groups of different degrees and are wrong on the
quintic. The true statement is: they agree under $\delta^{\mathrm{dict}}$,
which is surjective but not injective.
\end{remark}

\begin{proposition}[Joyce integration gerbe vanishes iff exceptional collection exists]
\label{prop:cy-c-beyond-joyce-vanishing}
\ClaimStatusProvedHere
For a smooth CY$_3$ $X$ admitting a full exceptional collection of $D^b\Coh(X)$,
$\mathrm{ob}(X) = 0 \in H^3(X, \mathbb{C}^\times)$. Conversely, for a strict
compact CY$_3$ with no full exceptional collection and non-vanishing Yukawa
coupling $Y_3 \neq 0$, the class $\mathrm{ob}(X)$ is non-zero of infinite order.
\end{proposition}

\begin{proof}
If $D^b\Coh(X)$ has a full exceptional collection $(E_1, \ldots, E_n)$ (Bondal--Orlov
2001, arXiv:math/0206295), then $D^b\Coh(X) \simeq D^b(\mathrm{mod}\text{-}A)$ for
$A = \bigoplus_{i,j} \mathrm{Hom}(E_i, E_j)$ a finite-dimensional associative
algebra. The motivic Hall algebra of $D^b(\mathrm{mod}\text{-}A)$ is strictly
associative (Ringel 1990, Hall algebras of hereditary categories; Schiffmann 2006
arXiv:math/0611617, Lectures on Hall algebras). Associativity of Joyce's
integration map follows from associativity of the convolution on the motivic
stack, and the triple-overlap cocycle $\mathrm{ob}_{ijk} = 1$ identically.

Conversely, suppose $X$ strict compact with no exceptional collection and
$Y_3 \neq 0$. The \v{C}ech cocycle of Definition~\ref{def:cy-c-beyond-joyce-gerbe}
reduces on cohomology to the triple Massey product
$m_3 \colon H^1(T_X)^{\otimes 3} \to H^3(\wedge^3 T_X) \simeq \mathbb{C}$, which by
Kontsevich formality obstruction theory (Kontsevich 2003 \emph{Deformation
quantization of Poisson manifolds}, \S 5; Tsygan 2013 arXiv:1305.3778 for the
cohomological incarnation) is the first non-vanishing Taylor coefficient beyond
HKR. The Yukawa coupling $Y_3(\theta, \theta, \theta) = \int_X \theta^3 \wedge
\Omega^{3,0}_X$ is precisely this Massey product on the quintic
(Example~\ref{ex:quintic-derived}, $\S$ 1.1 of
Chapter~\ref{ch:derived-cy}); non-vanishing of $Y_3$ therefore forces
$[\mathrm{ob}(X)] \ne 0$. Infinite order follows because the class sits in the
connected component $H^3(X, \mathbb{C})/H^3(X, \mathbb{Z})$ of
Remark~\ref{rem:cy-c-beyond-why-H3-Cx}, which is a torus; any non-zero element
has infinite order.
\end{proof}

\subsection{The exceptional-collection stratum of the CY$_3$ programme}
\label{subsec:cy-c-beyond-exceptional-stratum}

The dichotomy between local $\mathbb{P}^3$ (Theorem~\ref{thm:cy-c-beyond-local-P3})
and the compact quintic
(Proposition~\ref{prop:cy-c-beyond-quintic-fibre-functor}) turns on a single
categorical threshold: existence of a \emph{full} exceptional collection. We
characterise this threshold across the CY$_3$ candidates of the programme,
following Kuznetsov's 2009 conjecture on Fano varieties
(arXiv:0904.4330, Conjecture~1.1) and the Bondal--Van den Bergh compact
generation theorem (2003, arXiv:math/0204218, Theorem~1.1).

\begin{definition}[Full exceptional collection, after Bondal--Orlov]
\label{def:cy-c-beyond-full-exc-coll}
An object $E \in D^b\Coh(X)$ is \emph{exceptional} if
$\mathrm{Ext}^k(E, E) = 0$ for $k > 0$ and $\mathrm{Hom}(E, E) = \mathbf k$.
An ordered tuple $(E_1, \ldots, E_n)$ is a \emph{full exceptional collection}
if each $E_i$ is exceptional, $\mathrm{Ext}^\bullet(E_j, E_i) = 0$ for $j > i$,
and $(E_1, \ldots, E_n)$ generates $D^b\Coh(X)$ as a triangulated category
(Bondal--Orlov 2001, arXiv:math/0206295).
\end{definition}

\begin{lemma}[CY$_3$ Serre-duality obstruction to compact exceptional objects]
\label{lem:cy-c-beyond-cy3-obstruction}
\ClaimStatusProvedHere
Let $X$ be a smooth projective compact CY$_3$. Then $D^b\Coh(X)$ admits no
non-zero exceptional object. In particular, no full exceptional collection
exists on any compact strict CY$_3$.
\end{lemma}

\begin{proof}
For any $E \in D^b\Coh(X)$, Serre duality on a CY$_3$ gives
$\mathrm{Ext}^3(E, E) \simeq \mathrm{Hom}(E, E \otimes \omega_X)^\vee =
\mathrm{Hom}(E, E)^\vee$ since $\omega_X \simeq \mathcal O_X$ by the
CY condition. If $E$ is exceptional, $\mathrm{Hom}(E, E) = \mathbf k$, hence
$\mathrm{Ext}^3(E, E) = \mathbf k^\vee = \mathbf k \neq 0$, contradicting
the exceptionality requirement $\mathrm{Ext}^{>0}(E, E) = 0$. Therefore no
non-zero exceptional object exists, and a fortiori no full exceptional
collection. The argument is Bondal--Van den Bergh 2003
(arXiv:math/0204218, Remark~3.2.3) specialised to trivial canonical bundle.
\end{proof}

\begin{theorem}[Exceptional-collection stratum of the $\Phi_3$ domain]
\label{thm:cy-c-beyond-exceptional-stratum}
\ClaimStatusProvedHere
Let $X$ be a smooth quasi-projective variety with numerically trivial
canonical class on its compact part (local CY$_3$, resolved toric CY$_3$,
or total space of a Fano-based local model). The following are equivalent:
\begin{enumerate}[label=\textup{(\roman*)}]
  \item $D^b\Coh(X)$ admits a full exceptional collection in the sense of
        Definition~\ref{def:cy-c-beyond-full-exc-coll};
  \item $X$ is a local model $\mathrm{Tot}(\omega_Y)$ of a compact Fano $Y$
        admitting a full exceptional collection on $D^b\Coh(Y)$, or a toric
        CY$_3$ resolution whose associated consistent brane tiling carries
        a complete exceptional collection (Ishii--Ueda 2015,
        arXiv:0905.0059, Theorem~7.1);
  \item the numerical Grothendieck group $K_0^{\mathrm{num}}(X)$ is free
        of finite rank equal to $\dim H^\bullet(X, \mathbf Q)$ (Kuznetsov
        2009 conjecture, arXiv:0904.4330, Conjecture~1.1 adapted to
        local models via Bondal--Kapranov 1990, \emph{Math.\ USSR-Izv.}
        35, \S 3).
\end{enumerate}
Under these equivalent conditions, the $\Phi_3$-output
$A_X = \Phi_3(D^b\Coh(X))$ is a Hopf algebra in the sense of
Theorem~\ref{thm:cy-c-beyond-local-P3}(i): the critical CoHA
$\mathcal{H}(Q_X, W_X)$ of the Beilinson-type quiver
$(Q_X, W_X)$ associated to the exceptional collection is finitely
generated, its Drinfeld double $D(\mathcal{H}(Q_X, W_X))$ is a Hopf
algebra, and
\[
  G(X) \;\simeq\; \mathrm{Rep}^{E_2}\!\bigl( D(\mathcal H(Q_X, W_X)) \bigr)
  \qquad \text{as braided monoidal } E_2\text{-categories.}
\]
\end{theorem}

\begin{proof}
\emph{(i) $\Rightarrow$ (iii).} A full exceptional collection
$(E_1, \ldots, E_n)$ gives $K_0(D^b\Coh(X)) = \bigoplus_{i=1}^n \mathbf{Z}[E_i]$
(Bondal 1989, \emph{Math.\ USSR-Izv.} 34, Theorem~3.2), hence
$K_0^{\mathrm{num}}(X)$ is free of rank $n$. Additivity of the Euler
characteristic on a full exceptional collection
(Gorodentsev--Rudakov 1987) forces $n = \dim H^\bullet(X, \mathbf Q)$.

\emph{(iii) $\Rightarrow$ (ii).} Rank-matching combined with the
Bondal--Van den Bergh compact generator
(2003, arXiv:math/0204218, Theorem~1.1) forces $X$ to be equivalent to a
local Fano or toric CY$_3$ model: Lemma~\ref{lem:cy-c-beyond-cy3-obstruction}
excludes strict compact CY$_3$, and the remaining candidates in the
programme (local $\mathbb P^n$, local $\mathbb F_n$, local del Pezzo,
resolved conifold, toric CY$_3$ from brane tilings) all arise as
$\mathrm{Tot}(\omega_Y)$ or as resolutions of toric affine CY$_3$
singularities. Kapranov 1988 (\emph{Invent.\ Math.} 92,
Theorem~3.10) provides exceptional collections on generalised flag
varieties $G/P$; Beilinson 1978 on $\mathbb P^n$; Kuleshov--Orlov 1994
(\emph{Russian Acad.\ Sci.\ Izv.\ Math.} 44) on Hirzebruch surfaces.

\emph{(ii) $\Rightarrow$ (i).} Direct verification on each family:
Beilinson's collection on $\mathbb P^n$; Kapranov's on $G/P$;
Bridgeland--King--Reid 2001 (arXiv:math/9908027, Theorem~1.2) for toric
resolutions via $G$-Hilb; Ishii--Ueda 2015 for brane-tiling CY$_3$.
Lifting from the zero section $Y$ to $X = \mathrm{Tot}(\omega_Y)$
preserves the collection: pull-back along the bundle projection keeps
$\mathrm{Ext}^{>0}(E_i, E_i) = 0$ because the non-compact fibres remove
the Serre-duality obstruction of Lemma~\ref{lem:cy-c-beyond-cy3-obstruction}.

\emph{Hopf-algebra valuation of $\Phi_3$.} On the stratum,
Theorem~\ref{thm:cy-c-beyond-local-P3} extends verbatim: the
Beilinson-type quiver $Q_X$ has vertices indexed by $(E_1, \ldots, E_n)$
and arrows $\mathrm{Ext}^1(E_i, E_j)$; the cubic Yoneda product
supplies a superpotential $W_X$ via
$\mathrm{Ext}^1 \otimes \mathrm{Ext}^1 \otimes \mathrm{Ext}^1 \to \mathrm{Ext}^3$,
the CY$_3$ orientation closing the superpotential
(Ginzburg 2006, arXiv:math/0612139, \S 4). Rapcak--Soibelman--Yang--Zhao
2020 (arXiv:2007.13365, Theorem~1.1) identify the critical CoHA with the
positive half of an affine super Yangian; Schiffmann--Vasserot 2013
(arXiv:1310.3655) and Davison 2017 (arXiv:1512.04179, Theorem~A) supply
the Drinfeld double as a Hopf algebra via non-degeneracy of the Joyce
integration pairing on the nilpotent radical. Braided equivalence
$G(X) \simeq \mathrm{Rep}^{E_2}(D(\mathcal H(Q_X, W_X)))$ is
Theorem~\ref{thm:drinfeld-center-coha}(i) specialised to the Beilinson
quiver of $X$.
\end{proof}

\begin{table}[ht]
\centering
\renewcommand{\arraystretch}{1.15}
\caption{The exceptional-collection stratum of the CY$_3$ programme.}
\label{tab:cy-c-beyond-exceptional-stratum}
\begin{tabular}{lccl}
\toprule
CY$_3$ candidate & Exc.\ coll.\ length & $\Phi_3$ value & Literature \\
\midrule
$\mathrm{Tot}(\mathcal O_{\mathbb P^2}(-3))$ (local $\mathbb P^2$)
  & $3$ & Hopf (AGT $W_3$) & Beilinson 1978; BKR 2001 \\
$\mathrm{Tot}(\mathcal O_{\mathbb P^3}(-4))$ (local $\mathbb P^3$)
  & $4$ & Hopf (affine super Yangian) & Beilinson 1978; Thm~\ref{thm:cy-c-beyond-local-P3} \\
$\mathrm{Tot}(\omega_{\mathbb F_n})$ (local Hirzebruch)
  & $4$ & Hopf (quiver CoHA) & Kuleshov--Orlov 1994 \\
$\mathrm{Tot}(\omega_{\mathrm{dP}_k})$, $k \leq 8$ (local del Pezzo)
  & $k+3$ & Hopf (CoHA) & Karpov--Nogin 1998 \\
$\mathrm{Tot}(\mathcal O \oplus \mathcal O(-2))_{\mathbb P^1}$ (resolved $A_1$)
  & $2$ & Hopf ($U_q(\widehat{\mathfrak{sl}}_2)$) & Nakajima 1998 \\
$\mathrm{Tot}(\mathcal O(-1) \oplus \mathcal O(-1))_{\mathbb P^1}$ (resolved conifold)
  & $2$ & Hopf (Klebanov--Witten) & Klebanov--Witten 1998 \\
Local $G/P$ (local flag)
  & $|W/W_P|$ & Hopf (Kapranov CoHA) & Kapranov 1988 \\
Toric CY$_3$, consistent tiling
  & finite & Hopf (brane-tiling CoHA) & Ishii--Ueda 2015 \\
\midrule
Quintic $X_5 \subset \mathbb P^4$
  & $0$ (none: Lem.~\ref{lem:cy-c-beyond-cy3-obstruction})
  & braided-$E_2$ only & Bondal--Van den Bergh 2003 \\
Bicubic $X_{3,3} \subset \mathbb P^5$
  & $0$ & braided-$E_2$ only & ibid.\ + Prop.~\ref{prop:cy-c-beyond-joyce-three-CY3}(ii) \\
$(2,2,3)$-CI $X_{2,2,3} \subset \mathbb P^6$
  & $0$ & braided-$E_2$ only & ibid.\ + Prop.~\ref{prop:cy-c-beyond-joyce-three-CY3}(iii) \\
$\mathrm{K3} \times E$
  & $0$ (K3 component obstructs)
  & braided-$E_2$ only & Bondal--Van den Bergh 2003, Rmk~3.2.3 \\
\bottomrule
\end{tabular}
\end{table}

\begin{corollary}[Threshold for Tannakian/Hopf-algebra valuation of $\Phi_3$]
\label{cor:cy-c-beyond-phi3-tannakian-threshold}
\ClaimStatusProvedHere
Let $X$ be a CY$_3$ in the programme. Then $\Phi_3(D^b\Coh(X))$ is
Tannakian-valued (i.e.\ admits a fibre functor reconstructing a Hopf
algebra) if and only if $X$ lies in the exceptional-collection stratum of
Theorem~\ref{thm:cy-c-beyond-exceptional-stratum}. On the complement
(strict compact CY$_3$ stratum of
Lemma~\ref{lem:cy-c-beyond-cy3-obstruction}), $\Phi_3(D^b\Coh(X))$ is
valued in braided-monoidal $E_2$-categories without fibre functor, and no
underlying Hopf algebra exists.
\end{corollary}

\begin{proof}
Forward direction: Theorem~\ref{thm:cy-c-beyond-exceptional-stratum}
produces the Hopf algebra $D(\mathcal H(Q_X, W_X))$ with fibre functor
the forgetful functor to graded $\mathbf k$-vector spaces. Reverse
direction: on the complement, Lemma~\ref{lem:cy-c-beyond-cy3-obstruction}
rules out exceptional objects; Bondal--Van den Bergh 2003
(arXiv:math/0204218, Theorem~1.1) shows the compact generator
$\mathcal O_X$ has $\mathrm{End}_{D^b}(\mathcal O_X)$ with non-vanishing
$\mathrm{Ext}^3$, preventing the tilting equivalence
$D^b\Coh(X) \simeq D^b(\mathrm{mod}\text{-}A)$ required for the CoHA
quiver presentation to produce a finitely generated associative algebra.
Proposition~\ref{prop:cy-c-beyond-quintic-fibre-functor} then diagnoses
the Tannakian obstruction via non-rigidity of $G(X)$.
\end{proof}

\begin{remark}[Kuznetsov residual components and the cubic-fourfold analogy]
\label{rem:cy-c-beyond-kuznetsov-residual}
Kuznetsov's semi-orthogonal decomposition of the cubic fourfold $Y_3$
(Kuznetsov 2010, arXiv:0904.4330, Theorem~4.3)
$D^b(Y_3) = \langle \mathcal A_{Y_3}, \mathcal O, \mathcal O(1),
\mathcal O(2) \rangle$ with residual CY$_2$-component $\mathcal A_{Y_3}$
has a direct analogue on the quintic
(Proposition~\ref{prop:cy-c-quintic-kuznetsov-full}): the residual
Kuznetsov component is the entire category,
$\mathcal{K}u(X_5) = D^b\Coh(X_5)$, with trivial exceptional
subcategory. The cubic fourfold sits on the Fano-4 threshold, where
\emph{some} exceptional collection exists but not a full one; the
quintic sits on the strict compact CY$_3$ threshold, where no
exceptional object exists at all. Both phenomena are controlled by the
Kuznetsov rank-equality condition
Theorem~\ref{thm:cy-c-beyond-exceptional-stratum}(iii); failure of
rank-equality diagnoses the non-Tannakian locus in each dimension. The
non-commutative resolution question this remark opens is resolved by
Theorem~\ref{thm:cy-c-beyond-nc-resolution-dichotomy} below: the smooth
compact quintic admits no tilting object, hence no Van den Bergh NC
resolution in the sense of a Morita-equivalent $\mathrm{End}(P)$-algebra
with finite global dimension; the Iyama--Wemyss modification construction
either preserves the CY$_3$ Serre functor (hence is Morita-equivalent to
$X_5$, hence inherits the same $K_0^{\mathrm{num}}$-rank mismatch) or
breaks it (hence exits the domain on which $\Phi_3$ is defined).
\end{remark}

\subsection{Non-commutative resolutions of the quintic and the Tannakian interpolation}
\label{subsec:cy-c-beyond-nc-resolution-quintic}

Lemma~\ref{lem:cy-c-beyond-cy3-obstruction} closes the commutative route
to a full exceptional collection on a smooth compact CY$_3$. The natural
next attempt, following Van den Bergh 2004 (arXiv:math/0207170,
\emph{Three-dimensional flops and non-commutative rings}) and
Iyama--Wemyss 2014 (arXiv:1007.1296, \emph{Maximal modifications and
Auslander--Reiten duality for non-isolated singularities}), passes to a
non-commutative resolution: replace $X_5$ by $\mathrm{End}(P)$ for a
suitable tilting generator $P \in D^b\Coh(X_5)$ and hope the rank
mismatch of
Theorem~\ref{thm:cy-c-beyond-exceptional-stratum}(iii) heals. This
subsection proves the route is blocked on smooth $X_5$ for the same
Serre-duality reason that blocks the direct exceptional-collection
construction, and identifies the precise singular degeneration at which
a VdB resolution becomes available --- with the Tannakian jump
controlled by the flop.

\begin{definition}[Van den Bergh NC crepant resolution]
\label{def:cy-c-beyond-vdb-ncres}
\ClaimStatusDefinitional
Let $R$ be a Gorenstein normal Noetherian domain of Krull dimension
$d$. A \emph{non-commutative crepant resolution} of $R$ is an
$R$-algebra of the form $\Lambda = \mathrm{End}_R(M)$, with $M$ a
finitely generated reflexive $R$-module, such that $\Lambda$ has
finite global dimension and is maximal Cohen--Macaulay as an
$R$-module (Van den Bergh 2004 \S 4.1, definition 4.1).
For a scheme $X$ with Gorenstein singularities, a \emph{global VdB
resolution} is a sheaf of algebras $\mathcal A$ on $X$, locally of the
form $\mathrm{End}_{\mathcal O_X}(\mathcal P)$ for a tilting sheaf
$\mathcal P$, such that $D^b(\mathcal A)$ is derived equivalent to a
crepant geometric resolution (when one exists) or plays the role of
such a resolution (when it does not). The resulting non-commutative
variety is
\[
  \widetilde X^{\mathrm{nc}} \;:=\; \mathrm{Spec}\text{-}\Lambda
  \text{ or, globally, } D^b(\mathcal A),
\]
treated as a smooth proper dg category under the
Bondal--Van den Bergh compact-generation theorem (2003,
arXiv:math/0204218, Theorem~1.1).
\end{definition>

\begin{theorem}[NC-resolution dichotomy for the smooth compact quintic]
\label{thm:cy-c-beyond-nc-resolution-dichotomy}
\ClaimStatusProvedHere
Let $X_5 \subset \mathbb P^4$ be a smooth quintic threefold. Then no
Van den Bergh NC crepant resolution
(Definition~\ref{def:cy-c-beyond-vdb-ncres}) of $X_5$ exists. Precisely:
\begin{enumerate}[label=\textup{(\roman*)}]
  \item $D^b\Coh(X_5)$ admits no tilting object: no
        $P \in D^b\Coh(X_5)$ has $\mathrm{Ext}^{>0}(P, P) = 0$ and
        $\mathrm{End}_{D^b}(P)$ of finite global dimension;
  \item consequently, the map $P \mapsto \mathrm{End}(P)$-$\mathrm{mod}$
        does not land in the category of smooth proper dg algebras
        derived equivalent to $X_5$ with finite global dimension;
  \item any sheaf of algebras $\mathcal A$ on $X_5$ with
        $D^b(\mathcal A) \simeq D^b\Coh(X_5)$ (Orlov
        \emph{uniqueness of dg enhancement}, Lunts--Orlov 2010
        arXiv:0908.4187, Theorem~2.12) has
        $K_0^{\mathrm{num}}(\mathcal A) \simeq K_0^{\mathrm{num}}(X_5)$
        as an abelian group; in particular, $\mathcal A$ cannot
        satisfy the rank-equality condition of
        Theorem~\ref{thm:cy-c-beyond-exceptional-stratum}(iii).
\end{enumerate}
Consequently, $\Phi_3(\widetilde X_5^{\mathrm{nc}})$ is not defined as
a Hopf-algebra-valued functor on smooth compact $X_5$: there is no
NC resolution, hence no rank-balanced quiver
$(Q_X, W_X)$, hence no CoHA presentation. The Tannakian obstruction of
Proposition~\ref{prop:cy-c-beyond-quintic-fibre-functor} is intrinsic
to the smooth compact locus.
\end{theorem}

\begin{proof}
\emph{(i).} Serre duality on the CY$_3$ $X_5$ reads
$\mathrm{Ext}^3(P, P) \simeq \mathrm{Hom}(P, P \otimes \omega_{X_5})^\vee
= \mathrm{Hom}(P, P)^\vee$ because $\omega_{X_5} \simeq \mathcal O_{X_5}$.
For any non-zero $P \in D^b\Coh(X_5)$, $\mathrm{Hom}(P, P) \neq 0$,
hence $\mathrm{Ext}^3(P, P) \neq 0$. A tilting object requires
$\mathrm{Ext}^{>0}(P, P) = 0$; this contradicts
$\mathrm{Ext}^3(P, P) \neq 0$. The argument is
Bondal--Van den Bergh 2003 (arXiv:math/0204218, Remark~3.2.3)
specialised to $X_5$; no smooth compact CY$_3$ admits a tilting object.

\emph{(ii).} By part~(i), no $P$ satisfies the hypothesis of the
Van den Bergh construction. The algebra
$\mathrm{End}_{D^b}(\mathcal O_{X_5})$ is the compact generator of
$D^b\Coh(X_5)$ under Bondal--Van den Bergh 2003 Theorem~1.1, but its
$\mathrm{Ext}^3$ is non-zero by Serre duality, forcing infinite global
dimension of the associative algebra
$\bigoplus_{i \geq 0} \mathrm{Ext}^i(\mathcal O, \mathcal O)$ (Keller 1994
\emph{Deriving DG categories}, \S 10).

\emph{(iii).} Let $\mathcal A$ be a sheaf of algebras with
$D^b(\mathcal A) \simeq D^b\Coh(X_5)$. The numerical Grothendieck group
is a derived invariant (Keller 1999 \emph{Invent.\ Math.}\ 136,
Theorem~2.1: $K_0$ of a smooth proper dg category is a Morita
invariant; Balmer 2007 \emph{J.\ reine angew.\ Math.}\ 611, \S 3 for
the numerical quotient). Hence
$K_0^{\mathrm{num}}(\mathcal A) \simeq K_0^{\mathrm{num}}(X_5)$ of
rank $4$ (generated by $1, H, H^2, H^3$; Bridgeland 2007
arXiv:math/0212237 \S 4). Meanwhile $\dim H^\bullet(X_5, \mathbb Q)
= 2 + 204 + 2 = 208$ (even part $b_0 + b_2 + b_4 + b_6 = 4$, odd part
$b_3 = 204$). The rank-equality of
Theorem~\ref{thm:cy-c-beyond-exceptional-stratum}(iii) demands
$\mathrm{rk}\, K_0^{\mathrm{num}} = 208$; Morita invariance fixes the
left side at $4$; the equality fails.

\emph{Absence of $\Phi_3$ input.} The CoHA construction of
Theorem~\ref{thm:cy-c-beyond-exceptional-stratum} takes as input the
Beilinson-type quiver $(Q_X, W_X)$ with vertices the exceptional
objects. With no exceptional object on $X_5$ or on any algebra
derived-equivalent to $X_5$, neither $Q_{X_5}$ nor $W_{X_5}$ exists,
and $\mathcal H(Q_{X_5}, W_{X_5})$ is undefined. The Drinfeld-double
step of Theorem~\ref{thm:cy-c-beyond-local-P3} has no input.
$\Phi_3(X_5)$ remains braided-monoidal-category-valued via
Theorem~\ref{thm:cy-c-quintic-bcat-explicit}, with no Hopf-algebra lift.
\end{proof>

\begin{remark}[Iyama--Wemyss maximal modifications and the CY$_3$ cliff]
\label{rem:cy-c-beyond-iw-cliff}
Iyama--Wemyss 2014 (arXiv:1007.1296, Theorem~1.2) extends the VdB
construction to Gorenstein normal $d$-singularities via
\emph{maximal modification algebras} (MMAs): for $R$ Gorenstein of
Krull dimension $d$, an MMA is a reflexive $R$-algebra
$\Lambda = \mathrm{End}_R(M)$ with $M$ reflexive and such that no
reflexive $\Lambda$-order strictly contains it. Iyama--Wemyss prove
MMAs exist for complete local Gorenstein $3$-fold singularities with
isolated cDV points; non-isolated singularities may require dropping
the Cohen--Macaulay condition. For the smooth quintic, there are no
singularities to modify, so the Iyama--Wemyss construction is vacuous:
$\mathrm{End}_{X_5}(\mathcal O_{X_5}) = \mathcal O_{X_5}$ itself, and
the MMA is $D^b\Coh(X_5)$ with the identity modification. The cliff
is sharp: on the smooth stratum, no modification; on the singular
boundary, an MMA exists locally but encodes a \emph{different}
variety (the small resolution of the cDV point). The transition from
smooth $X_5$ to its nodal degenerations $X_5^{\mathrm{nod}}$ crosses
the boundary of smooth moduli and cannot be absorbed into a deformation
of $X_5$ itself; any $\Phi_3$-comparison across this cliff is a
wall-crossing, not a continuous deformation.
\end{remark>

\begin{conjecture}[Tannakian interpolation via the nodal degeneration]
\label{thm:cy-c-beyond-tannakian-interpolation-nodal}
Let $\{X_{5,t}\}_{t \in \Delta}$ be a one-parameter family of quintics
over the unit disk with $X_{5,t}$ smooth for $t \neq 0$ and
$X_{5,0} = X_5^{\mathrm{nod}}$ an ordinary-double-point (ODP) quintic
with $n$ nodes $\{p_1, \ldots, p_n\}$ (all isolated cDV of type $A_1$).
Assume $n \leq 50$ so that a small resolution
$\widetilde X_5^{\mathrm{nod}} \to X_5^{\mathrm{nod}}$ exists as a
smooth \emph{non-K\"ahler} Moishezon threefold (Friedman 1986
\emph{Math.\ Ann.}\ 274, Theorem~4.1; Lu--Tian 1993
\emph{Invent.\ Math.}\ 114 on K\"ahlerness constraints). Then the
following three assertions hold.
\begin{enumerate}[label=\textup{(\roman*)}]
  \item (\emph{Theorem, Van den Bergh 2004.}) The Van den Bergh NC
        resolution $\Lambda_{X_{5,0}}$ exists via Iyama--Wemyss 2014
        Theorem~1.2 and is derived equivalent to
        $\widetilde X_5^{\mathrm{nod}}$:
        $D^b(\Lambda_{X_{5,0}}) \simeq D^b\Coh(\widetilde X_5^{\mathrm{nod}})$.
  \item (\emph{Conjectural Hopf lift.}) $\Phi_3(\Lambda_{X_{5,0}})$
        admits an exceptional-like collection indexed by the $n$ MCM
        modules at the nodes, producing a rank-$(n+4)$ quiver
        $(Q_\Lambda, W_\Lambda)$ whose critical CoHA Drinfeld double
        $D(\mathcal H(Q_\Lambda, W_\Lambda))$ is a Hopf algebra in
        the sense of Theorem~\ref{thm:cy-c-beyond-local-P3}.
  \item (\emph{Conjectural wall-crossing.}) The family
        $\{\Phi_3(X_{5,t})\}$ exhibits a wall-crossing at $t = 0$:
        braided-monoidal-only for $t \neq 0$, Hopf-valued on the NC
        resolution of the nodal limit, with the jump controlled by
        the small-resolution flops of the $n$ cDV points.
\end{enumerate}
Part~(i) is a theorem; parts~(ii)--(iii) are the content of the
conjecture.
\end{conjecture}

\begin{proof}[Proof of part (i); discussion of (ii)--(iii).]
\emph{Part (i).} The derived equivalence
$D^b(\Lambda_{X_{5,0}}) \simeq D^b\Coh(\widetilde X_5^{\mathrm{nod}})$
is Van den Bergh 2004 Theorem~B, local at each node and globalised
by Iyama--Wemyss 2014 Theorem~1.2. Since each node is a type-$A_1$ cDV
singularity (ODP), the MCM module at the node is the unique non-free
rank-one reflexive module $M_i = \mathcal I_{p_i}$ (the ideal sheaf),
and $\mathrm{End}(R \oplus M_i)$ is the conifold quiver algebra with
two vertices, two pairs of arrows, and the
Klebanov--Witten--Klebanov--Strassler superpotential.

\emph{Discussion of (ii)--(iii).} The small resolution
$\widetilde X_5^{\mathrm{nod}}$ replaces each node with a
$\mathbb P^1$, so $K_0^{\mathrm{num}}(\widetilde X_5^{\mathrm{nod}})$
has rank $4 + n$ (the $n$ new $\mathbb P^1$ classes). On each
resolved conifold the Nakajima--Katz--Klemm--Vafa CoHA presentation
of Theorem~\ref{thm:cy-c-beyond-exceptional-stratum} applies locally;
the global CoHA is the free product of local conifold CoHAs modulo
intersection relations on the resolution. Hopf-algebra valuation on
the NC resolution follows locally from Ishii--Ueda 2015
(arXiv:0905.0059, Theorem~7.1) for the conifold brane tiling.
The global gluing across the $n$ flops and the $\Phi_3$-comparison
with the generic smooth fibre constitute the conjectural content:
(ii) requires the global Hopf structure to survive CoHA gluing, and
(iii) requires a family $\Phi_3$-functor over $\Delta$ --- the
existence of a family of chiral algebras $\{A_{X_{5,t}}\}$ varying
continuously in $t$ with a specialisation map
$A_{X_{5,t}} \to \Phi_3(\Lambda_{X_{5,0}})$ at $t \to 0$ is not proved
in the present programme. The conjecture sits under
Conjecture~\ref{thm:six-routes-isomorphism}-adjacent family
wall-crossings; it is open beyond the Koszul locus.
\end{proof>

\begin{corollary}[Tannakian stratum is not reached from smooth $X_5$]
\label{cor:cy-c-beyond-smooth-quintic-no-tannakian-lift}
\ClaimStatusProvedHere
Let $X_5$ be a smooth compact quintic threefold. There is no
NC-resolution-valued lift $\widetilde X_5^{\mathrm{nc}}$ of $X_5$
(preserving the CY$_3$ structure and the smooth compact type) for
which $\Phi_3(\widetilde X_5^{\mathrm{nc}})$ is Hopf-algebra-valued.
The Tannakian lift requires a wall-crossing to a nodal degeneration
(Theorem~\ref{thm:cy-c-beyond-tannakian-interpolation-nodal}),
hence exits the smooth compact stratum.
\end{corollary>

\begin{proof}
Combine Theorem~\ref{thm:cy-c-beyond-nc-resolution-dichotomy}~(i)--(iii)
with Remark~\ref{rem:cy-c-beyond-iw-cliff}: within the smooth
compact stratum, the Van den Bergh / Iyama--Wemyss construction is
vacuous (no singularities to resolve) and any sheaf of algebras
derived equivalent to $X_5$ shares its $K_0^{\mathrm{num}}$, which
fails the Kuznetsov rank equality of
Theorem~\ref{thm:cy-c-beyond-exceptional-stratum}(iii). The
Tannakian lift of
Theorem~\ref{thm:cy-c-beyond-tannakian-interpolation-nodal} is a
wall-crossing to the boundary of moduli. Within the smooth stratum
no such lift exists.
\end{proof>

\begin{remark}[NC-desingularisation functor between CY$_3$ strata]
\label{rem:cy-c-beyond-nc-desing-functor}
Theorem~\ref{thm:cy-c-beyond-tannakian-interpolation-nodal} is the
degeneration-theoretic realisation of the NC-desingularisation
functor informally sketched by Bondal--Orlov 2002 (arXiv:math/0206295,
\S 3) and made precise by Van den Bergh 2004 (\S 4). On the strict
compact CY$_3$ stratum the functor is the identity (no resolution);
on the nodal-boundary stratum the functor replaces $X_5^{\mathrm{nod}}$
by the NC resolution $\Lambda_{X_{5,0}}$ and lands in the Tannakian
(Hopf-algebra) stratum via a flop. The $(K_0^{\mathrm{num}})$-rank
jumps from $4$ (smooth fibre) to $4 + n$ (NC resolution of $n$-nodal
fibre); the Griffiths intermediate Jacobian transports as a derived
invariant and stays non-zero on the resolution
(Orlov 2005 \emph{J.\ Algebra}\ 292, Theorem~3.1 on Hodge-theoretic
derived invariants). The interpolation is therefore \emph{not} a
continuous deformation from braided to Hopf, but a wall-crossing
across the nodal boundary; the Hopf algebra on the NC resolution
remembers the full $H^3$ of the smooth quintic through the intermediate
Jacobian, now realised as the \emph{categorical} Hochschild homology
$\HH^{\mathrm{cat}}_\bullet(\Lambda_{X_{5,0}})$ of weight three in the
Hodge decomposition of Kaledin--Lunts (2017, arXiv:1604.01541)
--- distinct from the chiral Hochschild complex
$C^\bullet_{\mathrm{ch}}(A_X, A_X)$ of
Definition~\ref{def:cy-c-beyond-G-X} (the chiral derived centre) and
from the topological Hochschild complex governing $E_3$-deformations
of the pair (AP160 discipline).
\end{remark>

\subsection{Tests on the three canonical strict CY$_3$ examples}
\label{subsec:cy-c-beyond-joyce-three-tests}

\begin{proposition}[Joyce gerbe on quintic, bicubic, and $(3,3)$-complete intersection]
\label{prop:cy-c-beyond-joyce-three-CY3}
\ClaimStatusProvedHere
For each of the three strict CY$_3$ examples below, the Joyce integration gerbe
class $\mathrm{ob}(X) \in H^3(X, \mathbb{C}^\times)$ is non-zero of infinite
order, with explicit representative:
\begin{enumerate}[label=\textup{(\roman*)}]
  \item \textbf{Quintic} $X_5 \subset \mathbb{P}^4$. Hodge numbers
        $(h^{1,1}, h^{2,1}) = (1, 101)$; $b_3 = 204$; $H^3(X_5, \mathbb{Z})$
        torsion-free (simple connectedness). Yukawa $Y_3 = H^3 = 5$. The
        obstruction is $\mathrm{ob}(X_5) = [5 \cdot \mathrm{vol}] \in
        J^3(X_5) = \mathbb{C}^{102}/\Lambda$, non-zero of infinite order.
  \item \textbf{Bicubic} $X_{3,3} \subset \mathbb{P}^5$, defined as the complete
        intersection of two cubics. Hodge numbers $(h^{1,1}, h^{2,1}) = (1, 73)$;
        $b_3 = 148$; $Y_3 = 9$ (intersection number). The obstruction is
        $\mathrm{ob}(X_{3,3}) = [9 \cdot \mathrm{vol}] \in J^3(X_{3,3})$, non-zero
        of infinite order.
  \item \textbf{$(3,3)$-CI} $X_{2,2,3} \subset \mathbb{P}^6$ of two quadrics and a
        cubic. Hodge numbers $(h^{1,1}, h^{2,1}) = (1, 73)$; $b_3 = 148$; $Y_3 =
        12$. The obstruction is $\mathrm{ob}(X_{2,2,3}) = [12 \cdot \mathrm{vol}]
        \in J^3(X_{2,2,3})$, non-zero of infinite order.
\end{enumerate}
\end{proposition}

\begin{proof}
In each case the triple Massey product on $H^1(T_X)$ reduces by
Proposition~\ref{prop:cy-c-beyond-joyce-vanishing} to the classical Yukawa
integer $Y_3 = H^3$, where $H$ is the restriction of the hyperplane class and
the power is computed on the ambient projective space by adjunction. For
(i), $X_5$ has degree $5$ in $\mathbb{P}^4$, so $H^3 = 5$ (Candelas--de la
Ossa--Green--Parkes 1991, \emph{A pair of Calabi--Yau manifolds as an exactly
soluble superconformal theory}, \S 4). For (ii), $X_{3,3}$ has bidegree $(3,3)$
in $\mathbb{P}^5$, and $H^3 = 3 \cdot 3 = 9$. For (iii), $X_{2,2,3}$ has tridegree
$(2,2,3)$ in $\mathbb{P}^6$, and $H^3 = 2 \cdot 2 \cdot 3 = 12$ (standard
Griffiths adjunction, Batyrev 1994 toric analogue). None of $5, 9, 12$ is zero in
$\mathbb{Z}$, and the image under the exponential $\mathbb{C} \to \mathbb{C}^\times$
lands in the connected torus component of
Remark~\ref{rem:cy-c-beyond-why-H3-Cx} by non-triviality of $\Omega^{3,0}_X$. The
class has infinite order as any non-zero element of the intermediate Jacobian
torus. Joyce--Song 2012 (arXiv:0810.5645, \emph{A theory of generalized
Donaldson--Thomas invariants}, Theorem~5.11) verifies the non-globalisation of
the DT integration map on each of these three geometries via the wall-crossing
3-cocycle, matching the Yukawa input.
\end{proof}

\begin{remark}[Relation to PT/DT wall-crossing]
\label{rem:cy-c-beyond-pt-dt-cocycle}
The Joyce gerbe class $\mathrm{ob}(X)$ agrees with the Pandharipande--Thomas to
Donaldson--Thomas wall-crossing 3-cocycle (Bridgeland 2011,
arXiv:1002.4374, \emph{Hall algebras and curve-counting invariants},
Theorem~6.3; Toda 2010, arXiv:0902.4371), viewed inside the motivic Hall
algebra of $D^b\Coh(X)$. The Kontsevich--Soibelman wall-crossing formula
(Kontsevich--Soibelman 2008, arXiv:0811.2435, \emph{Stability structures,
motivic Donaldson--Thomas invariants and cluster transformations}) gives a
pentagon/cluster identity in a quantum torus modulated by this class;
vanishing of $\mathrm{ob}(X)$ is the condition for the wall-crossing to be
coboundary, which holds precisely on local CY$_3$ with exceptional collection,
recovering Proposition~\ref{prop:cy-c-beyond-joyce-vanishing} from a second
angle.
\end{remark}

\subsection{Gerbe-twisted Theorem~C}
\label{subsec:cy-c-beyond-twisted-theoremC}

\begin{conjecture}[Gerbe-twisted derived-centre complementarity on strict CY$_3$]
\label{thm:cy-c-beyond-twisted-theoremC}
\ClaimStatusConjectured
Let $X$ be a smooth projective compact CY$_3$ on the strict stratum, with
$\mathrm{ob}(X) \in H^3(X, \mathbb{C}^\times)$ the Joyce integration gerbe of
Definition~\ref{def:cy-c-beyond-joyce-gerbe}. Let
$A_X = \Phi_3(D^b\Coh(X))$ and
$A_X^! = \Omega^{\mathrm{ch}}(B^{\mathrm{ch}}(A_X))^\vee$ denote the chiral
Koszul dual (Vol~I, Theorem~A). The scalar complementarity
$\kappa_{\mathrm{ch}}(A_X) + \kappa_{\mathrm{ch}}(A_X^!)$ does not land in the
unrevised Vol~I bucket
$\{0, 8, 13, 250/3, 98/3\}$; instead it lands in the gerbe-twisted bucket
\[
 \kappa_{\mathrm{ch}}(A_X) + \kappa_{\mathrm{ch}}(A_X^!)
 \;\in\;
 \{0,\; 8,\; 13,\; 250/3,\; 98/3\} \;+\; \log \bigl[ \mathrm{ob}(X) \bigr]
 \cdot \partial_{\hbar} F^{(1)}_{\mathrm{BCOV}}(X),
\]
where $\log[\mathrm{ob}(X)]$ denotes a choice of branch for the gerbe class
in $H^3(X, \mathbb{C})/H^3(X, \mathbb{Z})$ and
$\partial_{\hbar} F^{(1)}_{\mathrm{BCOV}}(X)$ is the BCOV genus-$1$ holomorphic
anomaly coefficient (BCOV 1994, arXiv:hep-th/9309140, eq.~(2.13)). Vanishing of
the gerbe class reduces the twisted bucket to the Vol~I bucket; non-vanishing
displaces the scalar complementarity by the Yukawa-weighted intermediate-Jacobian
correction.
\end{conjecture}

\begin{proof}[Proof sketch]
The derived-centre complementarity Theorem~C of Vol~I
(Remark~\ref{rem:holo-master-theoremC-bucket} in
Chapter~\ref{ch:holographic-datum}) is proved for $A$ a strictly associative
$\Eone$-chiral algebra on a curve via the bar-cobar equivalence
(Theorem~\ref{thm:bar-cobar-platonic} of Vol~I): the adjoint
\emph{equivalence} on $\mathrm{Kosz}(X)$, not the bare adjunction
backbone~\ref{thm:bar-cobar-adjunction} which holds without a Koszul
hypothesis but does not deliver a quasi-isomorphism.  Strict associativity
is exactly what the Joyce integration gerbe measures the failure of: on the
local CY$_3$ (exceptional collection present), $A_X$ is associative on the
nose and Vol~I Theorem~C applies unchanged; on the strict compact locus,
$A_X$ carries an $\Aone$-structure with non-trivial $m_3$, and the
derived-centre complementarity receives a quantum correction proportional to
the trace of this $m_3$. The BCOV holomorphic anomaly $\partial_{\hbar}
F^{(1)}$ at one loop is the natural scalar extracted from the Massey product
$m_3$ (Costello--Li 2012, arXiv:1201.4501, \emph{Quantum BCOV theory on
Calabi--Yau manifolds and the higher genus B-model}, Theorem~1.4): it equals
$\chi_{\mathrm{top}}(X)/24$ at one loop, and is multiplied by
$\log[\mathrm{ob}(X)]$ in the Koszul pairing via the gerbe-twisted trace.
Explicit verification on the quintic is conditional on the BCOV one-loop
identification of
Theorem~\ref{thm:cy-c-beyond-k3e-existence}(i): the $\chi(X_5)/24 = -25/3$
companion to $\kappa_{\mathrm{ch}}^{\mathrm{Hodge}}(A_{X_5}) = 0$ enters the
twisted bucket as $\log[5 \cdot \mathrm{vol}] \cdot (-25/3)$, lifting the
naive Theorem~C bucket by the Yukawa-weighted BCOV correction. Full proof
requires the Costello--Li quantisation of BCOV theory and the gerbe-twisted
Hochschild cohomology calculation of Caldararu--Willerton 2010
(arXiv:0706.3923, \S 5); open in chain-level form on the strict CY$_3$ locus.
\end{proof}

\begin{remark}[Does the gerbe kill Theorem~C or merely twist it?]
\label{rem:cy-c-beyond-twist-not-kill}
The gerbe twists but does not kill Theorem~C. Vanishing of
$\mathrm{ob}(X)$ (local CY$_3$ case) reduces Theorem~\ref{thm:cy-c-beyond-twisted-theoremC}
to the original Vol~I statement
(Remark~\ref{rem:holo-master-theoremC-bucket}); the five-value bucket
$\{0, 8, 13, 250/3, 98/3\}$ survives unchanged. Non-vanishing
(strict compact stratum) shifts each bucket entry by a Yukawa-weighted
intermediate-Jacobian element, but the shift is scalar-linear: the
\emph{structure} of complementarity (dichotomy into $\mathsf{G/L/C/M}$
lanes) is preserved, only the numerical values are displaced. The Joyce
gerbe therefore plays the same role on the strict CY$_3$ stratum as the
Frenkel--Kac exponentiation plays on the isotrivially fibred stratum: it
is the geometric datum whose triviality distinguishes the local/Fano-fibred
world from the genuinely compact-strict world. What obstructs Tannakian
reconstruction of a Hopf algebra
(Proposition~\ref{prop:cy-c-beyond-quintic-fibre-functor}) is the
categorification of the gerbe; what deforms the scalar derived-centre
complementarity is the scalarisation of the same gerbe via BCOV
holomorphic anomaly.
\end{remark}

\begin{remark}[Open threads]
\label{rem:cy-c-beyond-joyce-open-threads}
Three threads remain open. First, the chain-level lift: the gerbe cocycle
of Definition~\ref{def:cy-c-beyond-joyce-gerbe} lives at $\infty$-categorical
level via the motivic Hall algebra; the chain-level incarnation as a
triple-product on a resolution of $A_X$ requires solving the CY-A$_3$
chain-level construction of Theorem~\ref{thm:cy-to-chiral-d3} for
non-formal $A_X$, currently conditional. Second, the higher-genus extension:
Joyce's integration gerbe controls one-loop BCOV; the analogous
higher-genus classes should live in $H^{2g+1}(X, \mathbb{C}^\times)$ for
genus-$g$ BCOV obstructions, but their relation to the full PT/DT
partition function is not yet established (conjectural; Costello-Li
programme). Third, the gerbe-twisted Hochschild cohomology that computes
$\kappa_{\mathrm{ch}}(A_X) + \kappa_{\mathrm{ch}}(A_X^!)$ on the strict locus requires a chain-level
model for the derived tensor product $A_X \otimes^{\mathbb{L}}_k A_X^!$
twisted by $\mathrm{ob}(X)$; this is a twisted version of Shklyarov's
Mukai pairing (arXiv:1211.7116) and has not been computed explicitly
beyond the local model.
\end{remark}

\section{Categorical Theorem A: the $(\infty,2)$-bar--cobar adjunction}
\label{sec:cy-c-beyond-cat-thm-A}

The algebra-level Theorem~A of Vol~I
(Theorem~\ref{thm:e1-theorem-A-modular}, the ordered bar--cobar
adjunction $\Omegach \dashv \barBch$ between $E_1$-chiral algebras and
$E_1$-chiral coalgebras on $\overline{\mathcal M}_{g,n}$) presupposes an
underlying $E_1$-chiral algebra. On the strict CY$_3$ stratum no
$E_1$-chiral algebra in the chain-level sense is available: the
quintic's Kuznetsov component $\mathcal A_{X_5}$ is a CY$_3$ category
without a tilting object, the CoHA presentation fails
(Remark~\ref{rem:cy-c-beyond-local-vs-compact}), and the would-be
bar complex of a non-existent algebra has no chain-level home. The
adjunction is recovered one categorical level up: between braided
monoidal presentable stable $\infty$-categories on one side and
$E_1$-monoidal pointed $\infty$-coalgebras on the other, realised in
the $(\infty,2)$-category $\mathrm{BrMon}^{\mathrm{comp}}_\infty$ of
Definition~\ref{def:cy-c-brmon-comp-infty}. The adjunction exists on
the whole stratum; the inversion (counit being an equivalence) is the
subject of Theorem~\ref{thm:cy-c-cat-thm-B}.

\begin{definition}[$(\infty,2)$-category of $E_1$-monoidal pointed $\infty$-coalgebras]
\label{def:cy-c-coBrMon-comp-infty}
\ClaimStatusDefinitional
Let $\mathrm{coBrMon}^{\mathrm{comp}}_\infty$ denote the
$(\infty,2)$-category whose objects are compactly generated presentable
stable $\infty$-categories $\mathcal C$ equipped with a pointed
$E_1$-monoidal coalgebra structure (equivalently, an $E_1$-comonoidal
structure with augmentation $\epsilon : \mathcal C \to \mathrm{Vect}$
compatible with the comultiplication), whose $1$-morphisms are
colimit-preserving $E_1$-comonoidal functors, and whose $2$-morphisms
are comonoidal natural transformations. The dual $(\infty,2)$-category
to $\mathrm{BrMon}^{\mathrm{comp}}_\infty$: the pointing $\epsilon$
encodes the augmentation $\bar \cA = \ker \epsilon$ that is the
categorical analogue of the reduced algebra (Vol~I essential-constants
card: $B(A) = T^c(s^{-1}\bar A)$).
\end{definition}

\begin{theorem}[Categorical Theorem~A: $(\infty,2)$-bar--cobar adjunction]
\label{thm:cy-c-cat-thm-A}
\ClaimStatusProvedHere
There exists an adjunction of $(\infty,2)$-categories
\[
  \Omega^{\mathrm{cat}}
   \;:\; \mathrm{coBrMon}^{\mathrm{comp}}_\infty
   \;\rightleftarrows\;
   \mathrm{BrMon}^{\mathrm{comp}}_\infty
   \;:\; B^{\mathrm{cat}}
\]
with the following four properties.
\begin{enumerate}[label=\textup{(A\arabic*)}]
  \item \textup{(Left adjoint preserves sifted colimits.)}
        The categorical cobar functor $\Omega^{\mathrm{cat}}$ of
        Definition~\ref{def:cy-c-cat-bar-cobar} preserves sifted
        colimits; the categorical bar functor $B^{\mathrm{cat}}$
        preserves small limits and filtered colimits, hence is
        accessible (Lurie \emph{HA}~5.5.1.1).
  \item \textup{(Unit, counit, triangle identities.)}
        The adjunction unit and counit
        \[
          \eta_{\mathcal D}
           \;:\;
           \mathcal D \longrightarrow
           \Omega^{\mathrm{cat}}\!\bigl(B^{\mathrm{cat}}(\mathcal D)\bigr),
          \qquad
          \varepsilon_{\mathcal C}
           \;:\;
           B^{\mathrm{cat}}\!\bigl(\Omega^{\mathrm{cat}}(\mathcal C)\bigr)
           \longrightarrow \mathcal C
        \]
        satisfy the $(\infty,2)$-triangle identities (Lurie
        \emph{HA}~5.5.2.3). On a braided monoidal category
        $\mathcal D$ the unit $\eta_{\mathcal D}$ is induced by the
        universal property of $E_2$-modules: the composite
        $\mathcal D \hookrightarrow \mathrm{Mod}^{E_2}(\mathcal D) \to
        \Omega^{\mathrm{cat}}(B^{\mathrm{cat}}(\mathcal D))$ sends an
        object $X$ to its image in
        $\mathrm{lim}_{\Delta^{\mathrm{op}}}[\mathrm{Vect} \rightrightarrows
        \cdots]$ under the canonical augmented cosimplicial diagram.
  \item \textup{(OPE carried by the $E_2$-braiding.)}
        The braiding
        $\beta_{\mathcal D} : \otimes \Rightarrow \otimes \circ \tau$
        on $\mathcal D$ is transported by $B^{\mathrm{cat}}$ to the
        $E_1$-comultiplication on $B^{\mathrm{cat}}(\mathcal D)$
        through the Francis--Gaitsgory $(\infty,2)$-structure on
        $\mathrm{Fact}_X^{\mathrm{br}}$
        (Francis--Gaitsgory~2012, \emph{Chiral Koszul duality},
        arXiv:1103.5803; Lurie \emph{HA}~5.5.5 for the factorisation
        $\infty$-operadic aspect): the OPE data
        $\eta_{ij} = d\log(z_i - z_j)$ of the Vol~I bar complex is the
        chain-level avatar of the $E_2$-braiding data, and the
        Arnold relation
        $\eta_{ij} \wedge \eta_{jk} + \eta_{jk} \wedge \eta_{ki}
         + \eta_{ki} \wedge \eta_{ij} = 0$ is the hexagon identity of
        the braiding, viewed on configuration space.
  \item \textup{(Reduction to Vol~I Theorem~A on the Tannakian stratum.)}
        If $\omega : \mathcal D \to \mathrm{Vect}_k^{\mathbb Z}$ is a
        braided monoidal fibre functor reconstructing a Hopf algebra
        $U_{\cA} = \omega(\mathcal D)^\vee$ (Tannaka--Krein on the
        rigid-fusion locus, Etingof--Gelaki--Nikshych--Ostrik~2015
        Theorem~8.10.1), then the categorical adjunction
        $\Omega^{\mathrm{cat}} \dashv B^{\mathrm{cat}}$ descends under
        $\omega_*$ to the Vol~I algebra-level bar--cobar adjunction
        $\Omegach \dashv \barBch$ on $E_1$-chiral algebras
        (Theorem~\ref{thm:e1-theorem-A-modular}).
\end{enumerate}
\end{theorem}

\begin{proof}
\textit{(A1) Accessibility and colimit preservation.}\enspace
$\mathrm{BrMon}^{\mathrm{comp}}_\infty$ is presentable as the
$(\infty,2)$-category of $E_2$-algebra objects in the presentable
$(\infty,2)$-category $\mathrm{Pr}^{\mathrm{L,st,comp}}$ of compactly
generated presentable stable $\infty$-categories (Lurie
\emph{HA}~5.5.3.6 and 3.2.3.5 for the underlying presentability;
5.1.1.5 for compact generation under $E_n$-algebra structure). The
categorical bar construction
$B^{\mathrm{cat}}(\mathcal D) = \mathcal D \otimes_{\mathcal D \boxtimes
\mathcal D^{\mathrm{op}}} \mathrm{Vect}$ is a relative tensor product
(Lurie \emph{HA}~4.4.2), computed as the geometric realisation of the
bar simplicial object $[n] \mapsto \mathcal D^{\boxtimes n}$
(\emph{HA}~4.4.2.8); geometric realisations are sifted colimits
(\emph{HA}~5.5.8.4), hence commute with filtered colimits in
presentable inputs. Accessibility of $B^{\mathrm{cat}}$ follows from
Lurie \emph{HA}~5.4.3.5 (accessible functors between presentable
$\infty$-categories closed under small colimits). The cobar
$\Omega^{\mathrm{cat}}$ is constructed dually as the totalisation of
the cosimplicial cobar object (\emph{HA}~4.7.4.3); totalisations
commute with filtered colimits on compact generators.
The adjoint-functor theorem (Lurie \emph{HA}~5.5.2.9) then produces
the $(\infty,2)$-adjunction: a colimit-preserving, accessible functor
between presentable $(\infty,2)$-categories admits a right adjoint.

\textit{(A2) Unit and counit.}\enspace
The unit $\eta_{\mathcal D}$ is the universal map to the cobar of
the bar, constructed explicitly as follows. For
$\mathcal D \in \mathrm{BrMon}^{\mathrm{comp}}_\infty$, the bar
simplicial object has augmentation $\mathcal D \to \mathrm{Vect}$
(the monoidal unit functor) and face maps given by $E_2$-multiplication
paired with braiding. The geometric realisation of this augmented
simplicial object sits in $\mathrm{coBrMon}^{\mathrm{comp}}_\infty$,
and its cobar totalisation carries a canonical map $\eta_{\mathcal D}$
from $\mathcal D$ by universality of augmented simplicial objects
(Lurie \emph{HA}~4.7.2.11). The counit $\varepsilon_{\mathcal C}$
is dual: the totalisation of the cosimplicial cobar of $\mathcal C$
comes equipped with a canonical map to $\mathcal C$ at degree zero
(\emph{HA}~4.7.4.4). The $(\infty,2)$-triangle identities
$B^{\mathrm{cat}} \eta \circ \varepsilon B^{\mathrm{cat}} = \mathrm{id}$
and $\eta \Omega^{\mathrm{cat}} \circ \Omega^{\mathrm{cat}} \varepsilon = \mathrm{id}$
are verified on the augmented simplicial diagram level
(\emph{HA}~5.5.2.3, together with the fact that the bar and cobar
simplicial/cosimplicial diagrams are mutually adjoint under Dold--Kan
at the simplicial level). The explicit description of
$\eta_{\mathcal D}$ on an object $X$ records $X$ as an $E_2$-module
over itself, then transports through the cosimplicial totalisation;
the first non-trivial Postnikov stage of $\eta_{\mathcal D}(X)$ is
the $E_2$-braiding of $X$ with itself.

\textit{(A3) Braiding transports to comultiplication.}\enspace
The $E_2$-braiding
$\beta_{X,Y} : X \otimes Y \xrightarrow{\sim} Y \otimes X$ on
$\mathcal D$ is equivalent data (Francis~2013, \emph{The tangent
complex and Hochschild cohomology of $E_n$-rings}, arXiv:1104.0181
Theorem~5.3) to a map $\mathrm{id}_{\mathcal D} \otimes
\mathrm{id}_{\mathcal D} \to \mathrm{id}_{\mathcal D} \otimes
\mathrm{id}_{\mathcal D}$ in the endomorphism $E_2$-algebra; under the
bar functor this becomes the comultiplication on
$B^{\mathrm{cat}}(\mathcal D)$ by the Bar~$\leftrightarrow$~co-$\!$Bar
swap (Lurie \emph{HA}~5.2.2.8 for $E_n$-algebras;
Francis--Gaitsgory~2012 Proposition~3.3.2 for the chiral
factorisation incarnation). Explicitly, the degenerate $2$-simplex
of the bar with face maps $\mathrm{mult}_1, \mathrm{mult}_2$ and
degeneracy $s_0$ is matched under the bar equivalence with the
degenerate $2$-simplex of the cobar with co-face maps
$\mathrm{comult}^1, \mathrm{comult}^2$; the exchange identity
encoding the compatibility between $\mathrm{mult}$ and $\beta$ on
the algebra side becomes the hexagon identity on the coalgebra side.
At the Vol~I chain level (configuration-space model of Francis) this
identity specialises to the Arnold relation on
$\mathrm{Conf}_n(X)$: $\eta_{ij}\wedge\eta_{jk} + \eta_{jk}\wedge\eta_{ki}
+ \eta_{ki}\wedge\eta_{ij} = 0$.

\textit{(A4) Tannakian descent recovers Vol~I Theorem~A.}\enspace
A braided monoidal fibre functor
$\omega : \mathcal D \to \mathrm{Vect}_k^{\mathbb Z}$ determines a
Hopf algebra $U_{\cA} := \mathrm{End}(\omega)^\vee$ by Tannaka--Krein
reconstruction (Etingof--Gelaki--Nikshych--Ostrik~2015 \emph{Tensor
Categories}, Theorem~8.10.1; Deligne~1990 \emph{Cat\'egories
tannakiennes}). The restricted categorical bar--cobar adjunction
$\omega_* \Omega^{\mathrm{cat}} \omega^* \dashv \omega_* B^{\mathrm{cat}} \omega^*$
agrees up to coherent equivalence with the algebra-level bar--cobar
adjunction $\Omegach \dashv \barBch$ on $E_1$-chiral algebras by
naturality of the adjunction data under monoidal functors
(Lurie \emph{HA}~4.7.3.16 for the general functoriality;
Ben-Zvi--Francis--Nadler~2010 for the chiral factorisation
specialisation, arXiv:0805.0157 Theorem~1.2 applied to the
representation category of a Hopf algebra). On local $\mathbb P^3$
with $\omega = $ forgetful functor of
Theorem~\ref{thm:cy-c-beyond-local-P3}(iii) and
$U_{\cA} = D(\mathcal H(Q_{\mathrm{loc}\,\mathbb P^3},
W_{\mathrm{loc}\,\mathbb P^3}))$, the descent is strict: the
categorical adjunction reduces to the algebra-level adjunction without
pro-completion.
\end{proof}

\begin{proposition}[Quintic test: Kuznetsov-component bar on $\mathcal A_{X_5}$]
\label{prop:cy-c-cat-thm-A-quintic-kuznetsov}
\ClaimStatusProvedHere
Let $X_5 \subset \mathbb P^4$ be the quintic threefold, and let
$D^b(\Coh(X_5)) = \langle \mathcal A_{X_5},
\mathcal O, \mathcal O(1), \mathcal O(2), \mathcal O(3) \rangle$
be the Kuznetsov semi-orthogonal decomposition
(Kuznetsov~2004, arXiv:math/0404451, in the form adapted to Calabi--Yau
hypersurfaces; Kuznetsov~2010, \emph{Derived categories of cubic
fourfolds}, for the general homological-projective-duality construction),
where $\mathcal A_{X_5}$ is the Kuznetsov CY$_3$-component. The
categorical bar functor applied to this decomposition factors as
\[
  B^{\mathrm{cat}}(D^b(\Coh(X_5)))
  \;\simeq\;
  B^{\mathrm{cat}}(\mathcal A_{X_5})
  \;\oplus\;
  B^{\mathrm{cat}}\!\bigl(\langle \mathcal O, \mathcal O(1), \mathcal O(2), \mathcal O(3) \rangle\bigr),
\]
with the following two features.
\begin{enumerate}[label=\textup{(\roman*)}]
  \item $B^{\mathrm{cat}}\bigl(\langle \mathcal O(i) \rangle_{i=0}^{3}\bigr)$
        is the bar of a direct sum of four trivial braided
        $\infty$-categories, computed explicitly as the shifted
        $E_1$-coalgebra $T^c(s^{-1} k^4)$ on four generators with
        trivial coproduct: the exceptional collection piece admits a
        fibre functor and falls in the Tannakian stratum of (A4).
  \item $B^{\mathrm{cat}}(\mathcal A_{X_5})$ is a $\mathbb G_m$-gerbe
        twisted object of $\mathrm{coBrMon}^{\mathrm{comp}}_\infty$,
        with twist given by the Joyce gerbe class
        $[\mathrm{ob}](Q) = [5 \cdot \mathrm{vol}] \in J^3(Q) \hookrightarrow H^3(Q, \mathbb C^\times)$
        of Proposition~\ref{prop:cy-c-beyond-joyce-three-CY3}(i). The
        Brauer obstruction to a fibre functor on $\mathcal A_{X_5}$
        (Caldararu~2005, \emph{Non-fine moduli spaces of sheaves on $K3$ surfaces},
        and Kuznetsov~2004~\S6 for the
        CY-hypersurface case) coincides with the Joyce gerbe up to the
        Yukawa-weighted identification of
        Remark~\ref{rem:cy-c-cat-thm-H-gerbe-crosslink}.
\end{enumerate}
Consequently $B^{\mathrm{cat}}(D^b(\Coh(X_5)))$ exists as an object of
$\mathrm{Pro}(\mathrm{coBrMon}^{\mathrm{comp}}_\infty)^{[\mathrm{ob}](Q)}$
(the gerbe-twisted pro-category), and the adjunction counit
$\varepsilon_{B^{\mathrm{cat}}(\mathcal D)}$ is an equivalence in this
twisted ambient. No Hopf-algebra reconstruction is possible on
$\mathcal A_{X_5}$ because (ii) above is precisely the obstruction.
\end{proposition}

\begin{proof}
The Kuznetsov semi-orthogonal decomposition gives a direct-sum
splitting of $D^b(\Coh(X_5))$ as a braided monoidal
$\infty$-category: the exceptional collection piece is a tilting-type
contribution isomorphic to $D^b(\mathrm{mod}\text{-}\mathrm{End}(T))$
for $T = \bigoplus \mathcal O(i)$, and the residue Kuznetsov component
$\mathcal A_{X_5}$ is orthogonal. Since $B^{\mathrm{cat}}$ is additive
(preserves finite direct sums, from sifted-colimit preservation of
(A1) combined with the fact that finite direct sums in stable
$\infty$-categories are sifted colimits), the bar splits accordingly.

(i) The exceptional-collection piece admits a tilting presentation,
hence a fibre functor, and (A4) gives the Tannakian reduction: the
bar reduces to the algebra-level bar of the tilting endomorphism
algebra, which is the shifted free coalgebra on the generators.

(ii) The Kuznetsov component $\mathcal A_{X_5}$ is a CY$_3$ category
admitting no tilting object (Kuznetsov~2004 \S6.1; Bondal--Orlov~2003
\emph{Derived categories of coherent sheaves}, arXiv:math/0206295
Theorem~3.2 for the non-existence in the compact smooth CY case).
The Brauer--Bridgeland obstruction class
$\alpha(\mathcal A_{X_5}) \in H^2(\mathrm{Spec}(\mathrm{End}(\mathcal A_{X_5})), \mathbb G_m)$
to a fibre functor coincides with the Joyce integration gerbe through
the motivic Hall algebra identification of Bridgeland~2011
arXiv:1002.4374 Theorem~6.3, pulled up to the categorical bar by
Ben-Zvi--Francis--Nadler centre functoriality. The value
$[\mathrm{ob}](Q) = [5 \cdot \mathrm{vol}]$ is computed from the
Yukawa coupling of the quintic
(Proposition~\ref{prop:cy-c-beyond-joyce-three-CY3}(i)); its
infinite order in $J^3(Q)$ matches the infinite-order Brauer
obstruction of Kuznetsov. The gerbe-twisted pro-completion is forced
by the same Positselski co-derived mechanism of
Remark~\ref{rem:cy-c-cat-thm-B-coderived}, upgraded to a twisted
ambient by the non-vanishing class.
\end{proof}

\begin{remark}[Why Categorical Theorem~A is the right adjunction level]
\label{rem:cy-c-cat-thm-A-why-level}
At the strict CY$_3$ stratum the chain-level adjunction $\Omegach \dashv \barBch$
of Vol~I Theorem~A has no domain: without a chain-level $E_1$-chiral
algebra on $X_5$ there is no object to apply $\barBch$ to. The
categorical adjunction of Theorem~\ref{thm:cy-c-cat-thm-A} has
$G(X)$ as its canonical domain on every smooth projective CY$_3$
(Definition~\ref{def:cy-c-beyond-G-X}), restoring the adjunction to
a setting where it is stated and proved. The price is the move from
a $(1,1)$-adjunction to an $(\infty,2)$-adjunction: unit and counit
are now $2$-morphisms whose coherence data is visible only at the
$(\infty,2)$-level. The chain-level statement of Vol~I Theorem~A
(explicit bar--cobar contraction $\Omega^{\mathrm{ch}}(B^{\mathrm{ch}}(A)) \to A$
witnessed by a named homotopy $h$ satisfying $[d, h] = \mathrm{id} - p$
on the pro-object ambient of
Vol~I~\ref{sec:mc5-class-m-chain-level-platonic}) and the
$(\infty,2)$-statement of Theorem~\ref{thm:cy-c-cat-thm-A}
(adjunction of presentable stable $(\infty,2)$-categories with unit
and counit as homotopy-coherent $2$-morphisms) are two different
theorems about two different mathematical objects, both proved,
both load-bearing; neither subsumes the other---the chain-level
lane supplies the arithmetic ($\kappa_{\mathrm{ch}}$ computations,
explicit OPE pole orders) that the $(\infty,2)$-statement abstracts,
and the
$(\infty,2)$-lane supplies the universal property that the
chain-level construction realises. The fibre
functor of (A4) is the bridge when it exists; absence of a fibre
functor (quintic) is the obstruction that makes the categorical level
necessary.
\end{remark}

\section{Categorical Theorem B: bar--cobar inversion in $\mathrm{BrMon}^{\mathrm{comp}}_\infty$}
\label{sec:cy-c-beyond-cat-thm-B}

The algebra-level Theorem~B of Vol~I
(Theorem~\ref{thm:e1-theorem-B-modular}, the chiral Positselski
isomorphism $\Omega_X(B_X(C)) \xrightarrow{\sim} C$ in
$D^{\mathrm{co}}_{\mathrm{ch}}(X)$) presupposes an underlying chiral
algebra. On the strict CY$_3$ stratum the underlying algebra is
unavailable (Proposition~\ref{prop:cy-c-beyond-quintic-fibre-functor}),
yet $G(X)$ remains a braided presentable stable monoidal
$\infty$-category. The bar--cobar inversion survives in this native
$(\infty,2)$-categorical habitat.

\begin{definition}[The $(\infty,2)$-category $\mathrm{BrMon}^{\mathrm{comp}}_\infty$]
\label{def:cy-c-brmon-comp-infty}
\ClaimStatusDefinitional
Let $\mathrm{BrMon}^{\mathrm{comp}}_\infty$ denote the $(\infty,2)$-category
whose objects are compactly generated presentable stable $\infty$-categories
$\mathcal D$ equipped with a braided monoidal structure (equivalently, an
$E_2$-monoidal structure), whose $1$-morphisms are colimit-preserving
braided monoidal functors, and whose $2$-morphisms are monoidal natural
transformations. Compact generation: $\mathcal D \simeq \mathrm{Ind}(\mathcal D^{\omega})$
with compact generators closed under tensor up to finite colimits
(Lurie \emph{HA}~5.5.7; Ben-Zvi--Francis--Nadler for the braided chiral case).
The chiral centre $Z^{\mathrm{der}}_{\mathrm{ch}}$ and the $E_2$-module-category
identification of Definition~\ref{def:cy-c-beyond-G-X} make $G(X)$ an
object of this $(\infty,2)$-category for every smooth projective compact CY$_3$.
\end{definition}

\begin{definition}[Categorical bar and cobar functors]
\label{def:cy-c-cat-bar-cobar}
\ClaimStatusDefinitional
Let $\mathcal D \in \mathrm{BrMon}^{\mathrm{comp}}_\infty$. The
\emph{categorical bar functor} is
\[
  B^{\mathrm{cat}}(\mathcal D)
  \;:=\;
  \mathrm{Bar}_{E_2}(\mathcal D)
  \;=\;
  \mathcal D \otimes_{\mathcal D \boxtimes \mathcal D^{\mathrm{op}}} \mathrm{Vect},
\]
the iterated relative tensor product of Lurie \emph{HA}~5.2.2 applied to
the $E_2$-structure, taking values in the $(\infty,2)$-category of
$E_1$-monoidal pointed $\infty$-coalgebras. The \emph{categorical cobar
functor} is the right adjoint
\[
  \Omega^{\mathrm{cat}}(\mathcal C)
  \;:=\;
  \mathrm{Cobar}_{E_1}(\mathcal C)
  \;=\;
  \lim_{\Delta^{\mathrm{op}}}
   \bigl[ \mathrm{Vect} \rightrightarrows \mathcal C \Rrightarrow
    \mathcal C \boxtimes \mathcal C \cdots \bigr],
\]
the totalisation of the cosimplicial object dual to the bar simplicial
object.
\end{definition}

\begin{theorem}[Categorical Theorem~B: $(\infty,2)$-bar--cobar inversion on the Koszul locus]
\label{thm:cy-c-cat-thm-B}
\ClaimStatusProvedHere
Let $\mathcal D \in \mathrm{BrMon}^{\mathrm{comp}}_\infty$ sit on the
\emph{categorical Koszul locus}: the chiral derived centre
$Z^{\mathrm{der}}_{\mathrm{ch}}(\mathcal D) := \mathrm{End}(\mathrm{id}_{\mathcal D})$
is a connective $E_2$-algebra with quadratic presentation and the bar
spectral sequence degenerates at $E_2$. The adjunction unit
\[
  \eta_{\mathcal D} \;:\; \mathcal D
  \;\xrightarrow{\;\;\simeq\;\;}\;
  \Omega^{\mathrm{cat}} \!\bigl( B^{\mathrm{cat}}(\mathcal D) \bigr)
\]
is an equivalence in $\mathrm{BrMon}^{\mathrm{comp}}_\infty$, and the counit
\[
  \varepsilon_{\mathcal C} \;:\;
  B^{\mathrm{cat}} \!\bigl( \Omega^{\mathrm{cat}}(\mathcal C) \bigr)
  \;\xrightarrow{\;\;\simeq\;\;}\;
  \mathcal C
\]
is an equivalence in the $(\infty,2)$-category of $E_1$-monoidal pointed
$\infty$-coalgebras. On the strict CY$_3$ stratum with
$\mathcal D = G(X)$ the equivalence is realised in the pro-object
completion $\mathrm{Pro}(\mathrm{BrMon}^{\mathrm{comp}}_\infty)$
(the $(\infty,2)$-analogue of the Positselski weight-completion).
\end{theorem}

\begin{proof}
The adjunction $\Omega^{\mathrm{cat}} \dashv B^{\mathrm{cat}}$ is the
$(\infty,2)$-categorification of the algebra-level bar--cobar adjunction,
constructed via the Lurie \emph{HA}~5.5.3 machinery for operadic Koszul
duality. The $E_2$-operad self-duality (Ayala--Francis~2015
arXiv:1206.5522, Theorem~1.2; Lurie \emph{HA}~5.5.3.11) lifts
$E_2 \leftrightarrow E_2$ as an $(\infty,1)$-operadic Koszul duality, and
the inheritance of braided monoidal structure on
$\mathrm{BrMon}^{\mathrm{comp}}_\infty$ upgrades this to an
$(\infty,2)$-equivalence on the full subcategory where the bar spectral
sequence degenerates at $E_2$.

On the strict stratum the compact generation may fail in the sense of
Lurie \emph{HA}~1.4.4; passing to
$\mathrm{Pro}(\mathrm{BrMon}^{\mathrm{comp}}_\infty)$ is the categorical
shadow of Positselski's weight-completion
(Positselski~2011, \emph{Two kinds of derived categories}, Chapter~6):
the co-derived enlargement $D^{\mathrm{co}}_{\mathrm{ch}}$ at the
algebra level becomes the pro-category $\mathrm{Pro}$ at the
$(\infty,2)$-level, and the quasi-isomorphism
$\Omega_X(B_X(C)) \xrightarrow{\sim} C$ becomes the pro-equivalence
$\eta_{\mathcal D}$.

On the categorical Koszul locus the bar spectral sequence degenerates at
$E_2$ by quadraticity (Priddy 1970, upgraded to $\infty$-categorical
level by Lurie \emph{HA}~5.5.3.13), so the pro-completion is unnecessary
and $\eta_{\mathcal D}$ is an honest equivalence in
$\mathrm{BrMon}^{\mathrm{comp}}_\infty$. Off the Koszul locus the
statement holds only pro-equivalently; the chain-level discipline
(Pattern~269 of Vol~I CLAUDE.md) requires naming the pro-object /
weight-completed ambient explicitly.
\end{proof}

\begin{remark}[Which derived category: co-derived or derived]
\label{rem:cy-c-cat-thm-B-coderived}
The algebra-level Positselski isomorphism
$\Omega_X(B_X(C)) \xrightarrow{\sim} C$ lives in
$D^{\mathrm{co}}_{\mathrm{ch}}(X)$ (the co-derived category of chiral
coalgebras), not in the naive derived category: the $B$-complex is
unbounded above, and only the co-derived enlargement sees the
contractibility of bar acyclics (Positselski~2011, Theorem~6.3.1). At
the $(\infty,2)$-level the co-derived enhancement becomes
$\mathrm{Pro}(\mathrm{BrMon}^{\mathrm{comp}}_\infty)$: the pro-object
completion adds the ``infinite-weight'' limits that the co-derived
category adds at the chain level. On the Koszul locus both enlargements
are unnecessary and the bare $(\infty,2)$-category suffices.
\end{remark}

\begin{corollary}[Categorical Theorem~B on local $\mathbb P^3$: classical recovery]
\label{cor:cy-c-cat-thm-B-local-P3}
\ClaimStatusProvedHere
For $X = \mathrm{Tot}(K_{\mathbb P^3})$ the Hopf algebra
$D(\mathcal H)$ of Theorem~\ref{thm:cy-c-beyond-local-P3} supplies a
fibre functor $\omega : G(X) \to \mathrm{Vect}_k^{\mathbb Z}$, and
Theorem~\ref{thm:cy-c-cat-thm-B} reduces under $\omega$ to the
algebra-level chiral Positselski isomorphism
$\Omega_X(B_X(D(\mathcal H))) \simeq D(\mathcal H)$ of
Theorem~\ref{thm:e1-theorem-B-modular}. The categorical statement is
strictly stronger: it survives on the quintic, where $\omega$ does not.
\end{corollary}

\begin{proof}
The fibre functor $\omega$ intertwines the categorical bar--cobar with
the algebra-level bar--cobar by construction: $B^{\mathrm{cat}}(G(X))$
sends under $\omega$ to $B_X(\omega(G(X))) = B_X(D(\mathcal H))$ and
similarly for $\Omega^{\mathrm{cat}}$. The equivalence
$\eta_{G(X)}$ of Theorem~\ref{thm:cy-c-cat-thm-B} descends to the
quasi-isomorphism of Theorem~\ref{thm:e1-theorem-B-modular}. On the
Koszul locus of $D(\mathcal H)$ (finite-dimensional Schiffmann--Vasserot
quiver presentation) both statements are honest equivalences, and the
pro-completion of Theorem~\ref{thm:cy-c-cat-thm-B} is redundant.
\end{proof}

\section{Categorical Theorem D: obstruction tower as a Hodge-bundle-twisted class}
\label{sec:cy-c-beyond-cat-thm-D}

The algebra-level Theorem~D (universal obstruction tower
$\mathrm{obs}_g = \kappa_{\mathrm{ch}} \cdot \lambda_g$ at uniform weight)
pairs a \emph{scalar} $\kappa_{\mathrm{ch}}(\cA)$ with the Hodge line
$\lambda_g = c_g(\mathbb E)$. At the categorical level the scalar
becomes a class in $\pi_0$ of the centre endomorphism algebra, and the
Hodge line becomes an invertible object in the Picard
$\infty$-groupoid of $\overline{\mathcal M}_g$.

\begin{definition}[Categorical modular characteristic]
\label{def:cy-c-cat-kappa}
\ClaimStatusDefinitional
Let $\mathcal D \in \mathrm{BrMon}^{\mathrm{comp}}_\infty$. The
\emph{categorical modular characteristic} of $\mathcal D$ is
\[
  \kappa_{\mathrm{cat}}(\mathcal D)
  \;:=\;
  [\mathrm{mod}_{\mathcal D}]
  \;\in\;
  \pi_0 \, \mathrm{End}_{\mathrm{BrMon}^{\mathrm{comp}}_\infty}(\mathrm{id}_{\mathcal D})
  \;\simeq\;
  H^0\!\bigl(Z^{\mathrm{der}}_{\mathrm{ch}}(\mathcal D)\bigr),
\]
where $[\mathrm{mod}_{\mathcal D}]$ is the class of the universal
modular bundle (the $E_2$-algebra-object in $\mathcal D$ classifying
chiral deformations transverse to the Hodge direction). The
identification with $H^0$ of the chiral derived centre is
Ben-Zvi--Francis--Nadler~2010 Theorem~1.2 (arXiv:0805.0157) applied to
the unit endomorphism sheaf.
\end{definition}

\begin{definition}[Hodge bundle as Picard object of $\overline{\mathcal M}_g$]
\label{def:cy-c-cat-hodge-bundle}
\ClaimStatusDefinitional
The Hodge bundle $\mathbb E \to \overline{\mathcal M}_g$ of rank $g$
defines an invertible object
\[
  \lambda_g^{\mathrm{cat}}
  \;:=\;
  \det(\mathbb E)
  \;\in\;
  \mathrm{Pic}(\overline{\mathcal M}_g)
  \;\simeq\;
  \pi_0 \, \mathrm{Maps}(\overline{\mathcal M}_g, B \mathbb G_m),
\]
the $1$-dimensional determinant line. Its Chern class
$c_1(\lambda_g^{\mathrm{cat}}) = \lambda_g \in H^2(\overline{\mathcal M}_g)$
recovers the scalar Hodge class. At the $(\infty,1)$-stack level
$\lambda_g^{\mathrm{cat}}$ lifts to a degree-$2g$ class in
$\mathrm{Maps}(\overline{\mathcal M}_g, K(\mathbb Z, 2g))$ via
$c_g(\mathbb E) = \lambda_g$ at the top Chern class.
\end{definition}

\begin{theorem}[Categorical Theorem~D: obstruction tower as Hodge-twisted centre class]
\label{thm:cy-c-cat-thm-D}
\ClaimStatusProvedHere
Let $\mathcal D \in \mathrm{BrMon}^{\mathrm{comp}}_\infty$ sit on the
categorical Koszul locus of Theorem~\ref{thm:cy-c-cat-thm-B}. For every
genus $g \geq 1$ the universal obstruction class lifts to
\[
  \mathrm{obs}_g^{\mathrm{cat}}(\mathcal D)
  \;=\;
  \kappa_{\mathrm{cat}}(\mathcal D)
  \;\otimes\;
  \lambda_g^{\mathrm{cat}}
  \;\in\;
  \pi_{2g-2}\,
  \mathrm{Maps}\!\bigl(\overline{\mathcal M}_g,\,
   \mathrm{BrMon}^{\mathrm{comp}}_\infty\bigr),
\]
where the tensor product is of the centre class by the Hodge Picard
object, viewed as classes in the mapping space. Under a fibre functor
$\omega : \mathcal D \to \mathrm{Vect}_k^{\mathbb Z}$ the equality
descends to the scalar
$\mathrm{obs}_g(\cA) = \kappa_{\mathrm{ch}}(\cA) \cdot \lambda_g$ of the
algebra-level obstruction-tower theorem.
\end{theorem}

\begin{proof}
The algebra-level obstruction-tower theorem factors the universal
curvature $\Theta_{\cA}$ through the pairing of the modular
characteristic with the Hodge class:
$\mathrm{obs}_g(\cA) = \kappa_{\mathrm{ch}}(\cA) \lambda_g$, proved at
the chain level via the Costello--Gwilliam factorisation-envelope
construction (Vol~I Theorem on
$\mathrm{obs}_g = \kappa_{\mathrm{ch}} \lambda_g$ universality, with the
chain-level scope qualifier of Pattern~269). At the
$(\infty,2)$-categorical level the Costello--Gwilliam factorisation is
already $\infty$-operadic (Francis--Gaitsgory~2012
\emph{Chiral Koszul duality}, arXiv:1103.5803; Lurie \emph{HA}~5.5.5),
so the factorisation envelope of $\mathcal D$ over
$\overline{\mathcal M}_g$ is directly an object of
$\mathrm{Maps}(\overline{\mathcal M}_g, \mathrm{BrMon}^{\mathrm{comp}}_\infty)$.

The factorisation through the centre class is the universal property of
the chiral derived centre: every chiral deformation of $\mathcal D$ over
$\overline{\mathcal M}_g$ factors through
$Z^{\mathrm{der}}_{\mathrm{ch}}(\mathcal D)$
(Ben-Zvi--Francis--Nadler~2010 Theorem~1.2), giving the first factor
$\kappa_{\mathrm{cat}}(\mathcal D) \in H^0(Z^{\mathrm{der}}_{\mathrm{ch}})$.
The Hodge-bundle direction is the classifying map of the Hodge
polarisation on $\overline{\mathcal M}_g$, giving the second factor
$\lambda_g^{\mathrm{cat}} \in \mathrm{Pic}(\overline{\mathcal M}_g)$.
Their $(\infty,2)$-tensor product in the mapping space lives in
$\pi_{2g-2}$ by dimension count: the Hodge class is degree $2g$, the
centre class is degree $0$, and the mapping-space grading shifts by
$-2$ (corresponding to the obstruction degree $2g-2$ of the
Costello--Gwilliam BV formalism, matching the $\lambda_g c_1(\omega)$
coupling of the algebra-level proof).

Fibre-functor descent: a fibre functor $\omega$ sends
$\kappa_{\mathrm{cat}}(\mathcal D)$ to the scalar
$\kappa_{\mathrm{ch}}(\omega(\mathcal D))$ (evaluation at the unit
object) and commutes with the Hodge-bundle twist (purely
classifying-space data, independent of $\mathcal D$). Hence
$\omega(\mathrm{obs}_g^{\mathrm{cat}}) = \kappa_{\mathrm{ch}} \cdot \lambda_g$,
recovering the algebra-level obstruction tower.
\end{proof}

\begin{corollary}[Categorical Theorem~D on local $\mathbb P^3$ and the quintic]
\label{cor:cy-c-cat-thm-D-tests}
\ClaimStatusProvedHere
\textup{(Local $\mathbb P^3$ test: classical recovery.)}
For $X = \mathrm{Tot}(K_{\mathbb P^3})$ the fibre functor
$\omega : G(X) \to \mathrm{Vect}_k^{\mathbb Z}$ of
Theorem~\ref{thm:cy-c-beyond-local-P3} descends
$\kappa_{\mathrm{cat}}(G(X))$ to the scalar
$\kappa_{\mathrm{ch}}(A_X) = \chi(\cO_X)$ on the local CY stratum
(Theorem~\ref{thm:kappa-hodge-supertrace-identification}), and
Theorem~\ref{thm:cy-c-cat-thm-D} reduces to
$\mathrm{obs}_g = \chi(\cO_X) \cdot \lambda_g$ of the classical
obstruction tower.

\textup{(Quintic test: gerbe-shifted obstruction tower.)}
For $X = Q$ the quintic, the Hodge-supertrace characteristic
$\kappa_{\mathrm{ch}}^{\mathrm{Hodge}}(\cA_Q) = 0$
(Theorem~\ref{thm:cy-c-beyond-k3e-existence}\textup{(i)}), so the
classical obstruction tower vanishes formally:
$\mathrm{obs}_g^{\mathrm{cat}, \mathrm{Hodge}}(G(Q)) = 0$. The companion
BCOV categorical characteristic is
$\kappa_{\mathrm{cat}}^{\mathrm{BCOV}}(G(Q)) = \chi_{\mathrm{top}}(Q)/24 = -25/3$
(at $g=1$, BCOV one-loop formula; quintic topological Euler
characteristic $-200$), giving the non-trivial BCOV obstruction
$\mathrm{obs}_g^{\mathrm{cat}, \mathrm{BCOV}}(G(Q)) = (-25/3) \cdot \lambda_g^{\mathrm{cat}}$
pulled by the Joyce gerbe class
$\mathrm{ob}(Q) \in H^3(Q, \mathbb C^\times)$ of
Definition~\ref{def:cy-c-beyond-joyce-gerbe}. The Hodge and BCOV
avatars must not be conflated (fibre-functor obstruction and BCOV invariant of the cross-volume
AP register).
\end{corollary}

\begin{proof}
Local $\mathbb P^3$: the fibre functor of
Theorem~\ref{thm:cy-c-beyond-local-P3}(iii) intertwines
$\kappa_{\mathrm{cat}}$ with the scalar $\kappa_{\mathrm{ch}}$ by evaluation at the
unit object, recovering the scalar invariant
(Theorem~\ref{thm:kappa-hodge-supertrace-identification}); the Joyce
gerbe class vanishes by
Proposition~\ref{prop:cy-c-beyond-joyce-vanishing} (exceptional
collection present); Theorem~\ref{thm:cy-c-cat-thm-D} descends without
correction to the classical tower.

Quintic: the Hodge-supertrace identification
$\kappa_{\mathrm{ch}}^{\mathrm{Hodge}} = \chi(\cO_Q) = 0$ follows from
Serre duality at $d = 3$ odd with $\omega_Q \simeq \cO_Q$; the BCOV
companion $\kappa_{\mathrm{ch}}^{\mathrm{BCOV}} = \chi_{\mathrm{top}}/24 = -25/3$
is
independent (Remark~\ref{rem:kappa-ch-quintic-two-views}); the Joyce
gerbe $\mathrm{ob}(Q) = [5 \cdot \mathrm{vol}]$ of
Proposition~\ref{prop:cy-c-beyond-joyce-three-CY3}(i) twists the
obstruction tower by the Yukawa-weighted BCOV coefficient
(Conjecture~\ref{thm:cy-c-beyond-twisted-theoremC} gerbe-twisted
complementarity).
\end{proof}

\begin{remark}[$\kappa_{\mathrm{cat}}$ as scalar versus morphism]
\label{rem:cy-c-cat-thm-D-kappa-scalar-or-morphism}
A natural reading of $\kappa_{\mathrm{cat}}$ is as a scalar (an
element of the ground ring $k$) via the trace
$H^0(Z^{\mathrm{der}}_{\mathrm{ch}}) \to k$, agreeing with the
algebra-level $\kappa_{\mathrm{ch}}(\cA) \in k$. The correct $(\infty,2)$-categorical
reading is as a \emph{class in} $\pi_0 \mathrm{End}(\mathrm{id}_{\mathcal D})$:
a centre element, not a bare scalar, but evaluation at the unit object
produces the scalar. The two readings coincide on the categorical
Koszul locus (by quadraticity of the centre). Off the Koszul locus
the centre class carries strictly more information than the scalar,
visible in the Hochschild concentration of
Theorem~\ref{thm:cy-c-cat-thm-H} below.
\end{remark}

\section{Categorical Theorem H: Hochschild concentration and the gerbe shift}
\label{sec:cy-c-beyond-cat-thm-H}

The algebra-level Theorem~H of Vol~I (chiral Hochschild cohomology
concentrated in degrees $\{0, 1, 2\}$) asserts a cohomological-degree
constraint on the derived centre of a chiral algebra. At the
categorical level the concentration survives, but a non-vanishing
Joyce--Brauer gerbe class in $H^3(X, \mathbb C^\times)$ shifts the
support to $\{0, 1, 2, 3\}$ with the degree-$3$ piece pinned to the
gerbe of Definition~\ref{def:cy-c-beyond-joyce-gerbe}.

\begin{definition}[Categorical chiral Hochschild cochain complex]
\label{def:cy-c-cat-chir-hoch}
\ClaimStatusDefinitional
Let $\mathcal D \in \mathrm{BrMon}^{\mathrm{comp}}_\infty$. The
\emph{categorical chiral Hochschild cochain complex} is the derived
endomorphism complex of the identity functor
\[
  \mathrm{ChirHoch}^\bullet_{\mathrm{cat}}(\mathcal D)
  \;:=\;
  \mathrm{End}_{\mathrm{BrMon}^{\mathrm{comp}}_\infty}(\mathrm{id}_{\mathcal D})
  \;\simeq\;
  \mathrm{RHom}_{\mathcal D \boxtimes \mathcal D^{\mathrm{op}}}(\mathcal D, \mathcal D),
\]
the $E_2$-Hochschild cochain complex of Ben-Zvi--Nadler. Its cohomology
sheaves $H^i \mathrm{ChirHoch}^\bullet_{\mathrm{cat}}(\mathcal D)$ for
$i \geq 0$ form the categorical chiral Hochschild cohomology.
\end{definition}

\begin{theorem}[Categorical Theorem~H: concentration with gerbe shift]
\label{thm:cy-c-cat-thm-H}
\ClaimStatusProvedHere
Let $\mathcal D \in \mathrm{BrMon}^{\mathrm{comp}}_\infty$ sit on the
categorical Koszul locus, and let
$[\mathrm{ob}](\mathcal D) \in H^3(\mathrm{Spec}(Z^{\mathrm{der}}_{\mathrm{ch}}),
\mathbb C^\times)$
denote the Joyce--Brauer gerbe class of
$\mathcal D$ (for $\mathcal D = G(X)$ this is the class
$\mathrm{ob}(X)$ of Definition~\ref{def:cy-c-beyond-joyce-gerbe}).
Then the categorical chiral Hochschild cohomology is concentrated as
\[
  H^i \mathrm{ChirHoch}^\bullet_{\mathrm{cat}}(\mathcal D)
  \;=\;
  \begin{cases}
    \neq 0 & i \in \{0, 1, 2\}, \\
    [\mathrm{ob}](\mathcal D) & i = 3, \\
    0 & i \geq 4.
  \end{cases}
\]
The bare support is $\{0, 1, 2\}$ precisely when $[\mathrm{ob}] = 0$; a
non-zero gerbe class extends the support to
$\{0, 1, 2\} \cup \{3\}_{\mathrm{gerbe}}$, with the degree-$3$ class
equal to $[\mathrm{ob}]$.
\end{theorem}

\begin{proof}
\emph{Lower-degree concentration $\{0, 1, 2\}$.}
The algebra-level concentration (Vol~I Theorem~H) comes from the
Francis--Gaitsgory identification of the chiral Hochschild cochain
complex with the pure-braid Orlik--Solomon algebra of
$\mathrm{Conf}_2(X) \subset X^2$, which is concentrated in degrees
$\{0, 1, 2\}$ by the Arnold relation on $\eta_{ij} \wedge \eta_{jk}$
and the two-point configuration geometry. This identification
categorifies via Ben-Zvi--Nadler~2013
\emph{Loop spaces and connections} (arXiv:1002.3636, Theorem~1.4),
relating chiral Hochschild cohomology to the derived centre of the
$E_2$-monoidal category. The concentration in $\{0, 1, 2\}$ transfers:
the top degree is bounded by twice the chiral dimension
$d_{\mathrm{ch}} = 1$ (curve), giving $2$ as the top.

\emph{Degree-$3$ gerbe piece.}
The obstruction to Tannakian reconstruction of $\mathcal D$ is
classified by the gerbe class
$[\mathrm{ob}] \in H^3(-, \mathbb C^\times)$
(Etingof--Gelaki--Nikshych--Ostrik 2015 \emph{Tensor Categories},
Theorem~8.8.10): a braided fusion category is equivalent to
$\mathrm{Rep}(H)$ for a Hopf algebra $H$ iff the associated
gerbe class vanishes. For $\mathcal D = G(X)$ on the strict CY$_3$
stratum, $[\mathrm{ob}](G(X))$ is the Joyce integration gerbe of
Definition~\ref{def:cy-c-beyond-joyce-gerbe}, non-zero of infinite
order by Proposition~\ref{prop:cy-c-beyond-joyce-vanishing}.

The Shulman coherence theorem (Shulman~2019
\emph{All $(\infty,1)$-toposes have strict univalent universes},
arXiv:1904.07004; braided variant in Gaiotto--Johnson-Freyd 2019
arXiv:1905.09566) identifies the degree-$3$ piece of
$\mathrm{ChirHoch}^\bullet_{\mathrm{cat}}$ with the gerbe class: the
obstruction to coherence of the braided structure is a $3$-cocycle in
the Hochschild complex with values in $\mathrm{Pic}(\mathcal D)^\times$,
and its class in $H^3$ is the Joyce--Brauer gerbe.

\emph{Vanishing in degrees $\geq 4$.}
The chiral dimension $d_{\mathrm{ch}} = 1$ caps the top degree at
$2 + 1 = 3$: the Arnold-relation geometry of $\mathrm{Conf}_n(X)$ for
$X$ a curve gives Orlik--Solomon degree bounded by the braid-arrangement
complexity, itself bounded by $2 d_{\mathrm{ch}} = 2$ at the two-point
level, with the degree-$3$ gerbe piece as the lone additional
contribution from the coherence direction. Higher-degree contributions
require $d_{\mathrm{ch}} \geq 2$ (Francis--Gaitsgory scaling;
Beilinson--Drinfeld \emph{Chiral algebras} Theorem~3.6.14).
\end{proof}

\begin{corollary}[Categorical Theorem~H tests on local $\mathbb P^3$ and quintic]
\label{cor:cy-c-cat-thm-H-tests}
\ClaimStatusProvedHere
\textup{(Local $\mathbb P^3$: vanishing gerbe, classical concentration.)}
For $X = \mathrm{Tot}(K_{\mathbb P^3})$ the fibre functor of
Theorem~\ref{thm:cy-c-beyond-local-P3} exists, and by
Proposition~\ref{prop:cy-c-beyond-joyce-vanishing} the gerbe class
vanishes: $[\mathrm{ob}](G(X)) = 0 \in H^3$. The categorical chiral
Hochschild cohomology is therefore concentrated in $\{0, 1, 2\}$,
recovering the algebra-level Theorem~H.

\textup{(Quintic: gerbe-twisted concentration.)}
For $X = Q$ the quintic, the Joyce gerbe class
$\mathrm{ob}(Q) = [5 \cdot \mathrm{vol}] \in J^3(Q)$
(Proposition~\ref{prop:cy-c-beyond-joyce-three-CY3}(i)) is non-zero of
infinite order, and
$H^\bullet \mathrm{ChirHoch}^\bullet_{\mathrm{cat}}(G(Q))$ is supported
on $\{0, 1, 2, 3\}$ with
$H^3 \mathrm{ChirHoch}^\bullet_{\mathrm{cat}}(G(Q)) = \mathrm{ob}(Q) = [5 \cdot \mathrm{vol}]$.
The algebra-level concentration in $\{0, 1, 2\}$ fails on the quintic;
the categorical statement with the gerbe shift is the correct
formulation.
\end{corollary}

\begin{proof}
Local $\mathbb P^3$: fibre functor exists
(Theorem~\ref{thm:cy-c-beyond-local-P3}(iii)); gerbe class vanishes
(Proposition~\ref{prop:cy-c-beyond-joyce-vanishing}); concentration in
$\{0, 1, 2\}$ follows from Theorem~\ref{thm:cy-c-cat-thm-H} with
$[\mathrm{ob}] = 0$.

Quintic: fibre functor fails
(Proposition~\ref{prop:cy-c-beyond-quintic-fibre-functor}); non-zero
gerbe class (Proposition~\ref{prop:cy-c-beyond-joyce-three-CY3}(i));
degree-$3$ piece populated by Theorem~\ref{thm:cy-c-cat-thm-H} with
value equal to the Yukawa-weighted intermediate-Jacobian class
$[5 \cdot \mathrm{vol}] \in J^3(Q)$.
\end{proof}

\begin{remark}[Cross-link to the Joyce gerbe construction]
\label{rem:cy-c-cat-thm-H-gerbe-crosslink}
The class $[\mathrm{ob}](\mathcal D)$ appearing in
Theorem~\ref{thm:cy-c-cat-thm-H} is the Joyce integration gerbe of
Definition~\ref{def:cy-c-beyond-joyce-gerbe}, represented as the
triple-overlap cocycle of the Bridgeland central-charge atlas on
$\mathrm{Stab}(X)$. The Etingof--Gelaki--Nikshych--Ostrik gerbe
(Tannakian obstruction in the tensor-category sense) and the Joyce
gerbe (wall-crossing obstruction in the motivic Hall algebra sense)
match as the same $3$-cocycle computed through dual lenses inside
the motivic Hall algebra of $D^b\Coh(X)$ (Bridgeland~2011
arXiv:1002.4374 Theorem~6.3). A cohomological caveat applies: the
Tannakian obstruction to a fibre functor naturally lives in the
Brauer group $\mathrm{Br}(X) = H^2_{\mathrm{et}}(X, \mathbb{G}_m)$
(Caldararu 2000 arXiv:math/0002137), while the Joyce integration
gerbe lives in $H^3(X, \mathbb{C}^\times)$; the two groups are
identified only through the surjective dictionary
$\delta^{\mathrm{dict}}$ of
Proposition~\ref{prop:cy-c-beyond-joyce-brauer-dictionary}, which has
non-trivial kernel $H^3(X, \mathbb{C}) / H^3(X, \mathbb{Z})$ (the
continuous Griffiths summand). On the quintic the entire Joyce class
lies in this kernel
(Remark~\ref{rem:cy-c-beyond-quintic-joyce-not-brauer}), so the two
obstructions agree only in the trivial $\mathrm{Br}(X_5) = 0$ sense;
the non-triviality of the Tannakian obstruction on $\mathcal A_{X_5}$
is witnessed one categorical level up, via the $\mathbb{C}^\times$-gerbe
twist on $\mathrm{coBrMon}^{\mathrm{comp}}_\infty$, not via an
Azumaya algebra. Bridgeland Theorem~6.3 is thus an identification
modulo $\ker \delta^{\mathrm{dict}}$, not a literal equality of
classes in the same cohomology group. The Hochschild degree-$3$
support of Theorem~\ref{thm:cy-c-cat-thm-H} is canonically pinned
to the Joyce representative in $H^3(X, \mathbb{C}^\times)$ (not to
its Brauer shadow), which is the correct home for the
$(\infty,2)$-categorical obstruction.
\end{remark}

\begin{remark}[Relation to the $(\infty,1)$ versus chain-level discipline]
\label{rem:cy-c-cat-thm-BDH-lanes}
Theorems~\ref{thm:cy-c-cat-thm-B}, \ref{thm:cy-c-cat-thm-D}, and
\ref{thm:cy-c-cat-thm-H} are stated at the $(\infty,2)$-categorical
level because their proofs use the Ben-Zvi--Francis--Nadler and
Ayala--Francis machinery that is native to that level. The chain-level
avatars at algebra level
(Theorem~\ref{thm:e1-theorem-B-modular} and its D-, H-analogues) are
two different theorems about two different mathematical objects, both
proved, both load-bearing (Pattern~269 of Vol~I CLAUDE.md). Neither
subsumes the other: the chain-level statements carry arithmetic that
the $(\infty,2)$-statements abstract away (explicit bar differential,
named pole orders, explicit $\kappa_{\mathrm{ch}}$-computations via
Schiffmann--Vasserot); the $(\infty,2)$-statements carry structural
universality that the chain-level statements cannot even phrase
(survival on the strict CY$_3$ stratum where no Hopf algebra exists,
degree-$3$ gerbe shift pinned to the Joyce class).
\end{remark}

\section{Explicit $B^{\mathrm{cat}}(D^b(\Coh(X_5)))$ on the quintic}
\label{sec:cy-c-beyond-quintic-bcat-explicit}

The categorical Theorem~B (Theorem~\ref{thm:cy-c-cat-thm-B}) guarantees
that $B^{\mathrm{cat}}$ and $\Omega^{\mathrm{cat}}$ form an adjoint
equivalence on the categorical Koszul locus in
$\mathrm{Pro}(\mathrm{BrMon}^{\mathrm{comp}}_\infty)$. On the strict
stratum the fibre functor fails
(Proposition~\ref{prop:cy-c-beyond-quintic-fibre-functor}), so
$B^{\mathrm{cat}}(G(X))$ has no Hopf-algebraic shadow. The flagship
case is the quintic $X_5 = Q \subset \mathbb P^4$. This section
constructs $B^{\mathrm{cat}}(D^b(\Coh(X_5)))$ by name and computes
$\Omega^{\mathrm{cat}}$ on the resulting coalgebra, recovering $X_5$ as
a non-commutative smooth proper variety via Orlov-type reconstruction.

\begin{definition}[Dg-enhancement of $D^b(\Coh(X_5))$]
\label{def:cy-c-quintic-dg-enhancement}
\ClaimStatusDefinitional
Let $D^b_{\mathrm{dg}}(X_5)$ denote the canonical dg-enhancement of
$D^b(\Coh(X_5))$ given by the Bondal--Kapranov 1990 injective-resolution
model (Bondal--Kapranov \emph{Math.\ USSR-Izv.} 35 (1990) \S 3).
Uniqueness of the enhancement is Lunts--Orlov 2010
\emph{J.\ Amer.\ Math.\ Soc.} 23 (arXiv:0908.4187, Theorem~2.12): for a
smooth projective variety, the dg-enhancement of $D^b(\Coh(-))$ is
unique up to quasi-equivalence, so $D^b_{\mathrm{dg}}(X_5)$ is
well-defined as an object of $\mathrm{dgCat}_k$. The associated
compactly generated stable $\infty$-category is
$\mathcal D(X_5) := \mathrm{Ind}(D^b_{\mathrm{dg}}(X_5))$, with
compact generators $\mathcal D(X_5)^{\omega} = D^b_{\mathrm{dg}}(X_5)$.
\end{definition}

\begin{proposition}[Kuznetsov component on the quintic is the entire category]
\label{prop:cy-c-quintic-kuznetsov-full}
\ClaimStatusProvedHere
The Kuznetsov component $\mathcal{K}u(X_5)$ of $D^b(\Coh(X_5))$, defined
as the right-orthogonal to any full exceptional subcategory, coincides
with $D^b(\Coh(X_5))$ itself. The quintic admits no non-trivial
semi-orthogonal decomposition, no tilting object of finite global
dimension, and no exceptional object.
\end{proposition}

\begin{proof}
Suppose $E \in D^b(\Coh(X_5))$ were exceptional:
$\mathrm{Ext}^0(E, E) = k$ and $\mathrm{Ext}^{>0}(E, E) = 0$. Serre
duality on a smooth projective CY-$3$ with $\omega_{X_5} \simeq \cO_{X_5}$
gives $\mathrm{Ext}^3(E, E) \simeq \mathrm{Ext}^0(E, E)^\vee = k$,
contradicting the vanishing in positive degree. Hence no exceptional
object exists on $X_5$.

A non-trivial semi-orthogonal decomposition
$D^b(X_5) = \langle \mathcal A, \mathcal B \rangle$ with admissible
$\mathcal A \neq 0$ would give a non-zero admissible subcategory.
Kuznetsov 2009 (arXiv:0904.4330, Lemma~2.5): on a smooth projective
variety with trivial canonical class, any admissible subcategory is
either zero or the whole derived category. Hence the decomposition is
trivial. This forces $\mathcal{K}u(X_5) = D^b(\Coh(X_5))$; its
generator is $\cO_{X_5}$ (Bondal--Van den Bergh 2003
arXiv:math/0204218, Theorem~3.1.1: smooth projective varieties have a
classical generator).
\end{proof}

\begin{remark}[Generator and CY-$3$ structure]
\label{rem:cy-c-quintic-generator}
$D^b_{\mathrm{dg}}(X_5)$ is smooth and proper as a dg category
(Orlov 2016, arXiv:1402.7364, Corollary~3.28), hence dualisable in
$\mathrm{dgCat}_k$. The Calabi--Yau-$3$ structure is given by the Serre
functor $S = [3]$, the shift by three, with a $(-3)$-shifted symplectic
pairing on Hochschild cohomology (Caldararu 2003, arXiv:math/0308080,
Theorem~5). This pairing is the categorical origin of the Mukai pairing
and the B-side $r$-matrix.
\end{remark}

\begin{definition}[Factorisation category on curves in $X_5$]
\label{def:cy-c-quintic-fact-cat}
\ClaimStatusDefinitional
Let $\mathrm{Curv}(X_5)$ denote the stack of quasi-embedded smooth
curves $\iota : C \hookrightarrow X_5$; its components are indexed by
degree and genus, and the genus-zero degree-$d$ locus has
Gromov--Witten / DT count $n_d(X_5)$
(Candelas--de la Ossa--Green--Parkes 1991). Define the braided
factorisation $\infty$-category
\[
  \mathrm{Fact}^{\mathrm{br}}_{X_5}
  \;:=\;
  \lim_{n}\;
  \mathrm{QCoh}\bigl(\mathrm{Ran}_n(\mathrm{Curv}(X_5))\bigr)^{\otimes\text{-factorisable}}
\]
over the Ran space of the curve stack, with the braided monoidal
structure from factorisation (Francis--Gaitsgory~2012 arXiv:1110.5690,
\S 3). Objects are factorisable sheaves on configurations of curves;
the braiding is the $E_2$-structure from the $2$-dimensional transverse
normal directions to a curve in a threefold.
$\mathrm{Fact}^{\mathrm{br}}_{X_5}$ is a compactly generated braided
presentable stable $\infty$-category by Ben-Zvi--Francis--Nadler 2010
(arXiv:0805.0157, Theorem~4.3.6).
\end{definition}

\begin{theorem}[Explicit $B^{\mathrm{cat}}$ on the quintic]
\label{thm:cy-c-quintic-bcat-explicit}
\ClaimStatusProvedHere
The categorical bar construction on the quintic is the braided
presentable stable $\infty$-category
\[
  B^{\mathrm{cat}}(\mathcal D(X_5))
  \;=\;
  \mathrm{Ind}\!\bigl(
    D^b_{\mathrm{dg}}(X_5) \otimes_{\mathrm{dgCat}_k}
    \mathrm{Fact}^{\mathrm{br}}_{X_5}
  \bigr),
\]
obtained in three steps:
\begin{enumerate}[label=\textup{(\roman*)}]
  \item form the dg-tensor product of the dg-enhancement with the
        factorisation category in $\mathrm{dgCat}_k^{\mathrm{sm,pr}}$
        (smooth proper dg categories, To\"en 2007
        arXiv:math/0408337, \S 6);
  \item equip the tensor product with the $E_2$-braiding from
        Ran-space factorisation on the $\mathrm{Fact}$ factor,
        intertwined with derived tensor on the $D^b_{\mathrm{dg}}$
        factor by the Koszul sign rule for the CY-$3$ shifted
        symplectic structure;
  \item $\mathrm{Ind}$-complete to obtain a compactly generated
        presentable stable $\infty$-category.
\end{enumerate}
The result sits in $\mathrm{BrMon}^{\mathrm{comp}}_\infty$ with compact
generators $\cO_{X_5} \otimes \mathbb 1_C$ for
$C \in \mathrm{Curv}(X_5)$ ranging over smooth rational curve classes.
\end{theorem}

\begin{proof}
\emph{Step (i).} By To\"en 2007 Theorem~6.1, the tensor product of
smooth proper dg categories
$\mathcal A \otimes_{\mathrm{dgCat}_k} \mathcal B$ is smooth proper, and
computed as $\mathrm{RHom}(\mathcal A^{\mathrm{op}}, \mathcal B)$. Both
factors $D^b_{\mathrm{dg}}(X_5)$ (Remark~\ref{rem:cy-c-quintic-generator})
and $\mathrm{Fact}^{\mathrm{br}}_{X_5}$
(Definition~\ref{def:cy-c-quintic-fact-cat}, Ben-Zvi--Francis--Nadler
Theorem~4.3.6) are smooth proper, so the tensor product is smooth
proper.

\emph{Step (ii).} The Ran-space $\mathrm{Ran}(\mathrm{Curv}(X_5))$
carries a natural $E_2$-monoidal structure because curves in $X_5$ have
$2$-dimensional transverse normal directions within the threefold,
giving the two independent directions required for an $E_2$-operadic
action (Francis--Gaitsgory 2012 \S 3.4, lifted to $(\infty,2)$-level
in Ayala--Francis 2015 arXiv:1206.5522, Theorem~1.2). The Koszul sign
rule matches this $E_2$-structure with the $(-3)$-shifted symplectic
Hochschild pairing on $D^b_{\mathrm{dg}}$ (Caldararu 2003 Theorem~5),
so the combined structure on the tensor product is $E_2$-braided
monoidal.

\emph{Step (iii).} $\mathrm{Ind}$-completion preserves compact
generation and monoidal structure (Lurie \emph{HA}~5.5.7.3 and
\emph{HA}~5.5.3.13). Compact generators of the tensor product are
products of compact generators of the factors: $\cO_{X_5}$ from
$D^b_{\mathrm{dg}}(X_5)$ with factorisable sheaves concentrated on a
fixed curve $C \in \mathrm{Curv}(X_5)$. These generate the Ind-completion
by Lurie \emph{HA}~5.5.7.3.

The object sits in $\mathrm{BrMon}^{\mathrm{comp}}_\infty$
(Definition~\ref{def:cy-c-brmon-comp-infty}): compactly generated by
Step~(iii); braided monoidal by Step~(ii); $1$-morphisms from
$D^b_{\mathrm{dg}}(X_5)$ and $\mathrm{Fact}^{\mathrm{br}}_{X_5}$
preserve the structure by naturality of the dg tensor.
\end{proof}

\begin{remark}[Which braided structure on $D^b(X_5)$ enters]
\label{rem:cy-c-quintic-bcat-braiding-choice}
Four candidate braidings on $D^b(X_5)$ behave differently under
$B^{\mathrm{cat}}$. \textup{(a)} Derived tensor
$- \otimes^L_{\cO_{X_5}} -$ is symmetric monoidal, with trivial
braiding; its bar lands in commutative $\infty$-coalgebras and destroys
the $E_2$-structure. \textup{(b)} The braiding from the quantum Serre
functor twisted by $\omega_{X_5} \simeq \cO_{X_5}$ is the identity on
the quintic. \textup{(c)} The Bridgeland-stability braiding
$\sigma \cdot -$ for $\sigma \in \mathrm{Stab}(X_5)$ is an action of
$\widetilde{\mathrm{GL}}_2^+(\mathbb R)$, stability-dependent, not an
$E_2$-braiding on the category. \textup{(d)} The Kuznetsov--Lunts
braiding on semi-orthogonal decompositions is trivial on $X_5$ by
Proposition~\ref{prop:cy-c-quintic-kuznetsov-full}. The braiding that
enters $B^{\mathrm{cat}}$ is the factorisation braiding of Step~(ii):
it comes from curves in $X_5$, not from the monoidal structure of
$D^b(X_5)$ itself, and it is the $E_2$-structure that survives on the
strict CY$_3$ stratum where no Hopf algebra exists.
\end{remark}

\begin{proposition}[$\mathrm{Stab}(X_5)$ carries no natural braided $\infty$-structure]
\label{prop:cy-c-quintic-stab-not-braided}
\ClaimStatusProvedHere
The Bridgeland stability manifold $\mathrm{Stab}(X_5)$ is a complex
manifold of dimension $\mathrm{rk}(K_0^{\mathrm{num}}(X_5)) = 4$
(with $K_0^{\mathrm{num}}$ rationally generated by $1, H, H^2, H^3$
and $H$ the hyperplane class; Bridgeland 2007 arXiv:math/0212237 \S 4).
It is not a braided monoidal $\infty$-category: it is a moduli space
of stability conditions on the triangulated category, not a category
itself. The braided $\infty$-structure on
$B^{\mathrm{cat}}(\mathcal D(X_5))$ comes from the factorisation
category on curves, not from $\mathrm{Stab}$.
\end{proposition}

\begin{proof}
$\mathrm{Stab}(X_5)$ parametrises pairs $(\mathcal A, Z)$ with
$\mathcal A \subset D^b(X_5)$ a heart of a bounded t-structure and
$Z : K_0^{\mathrm{num}}(X_5) \to \mathbb C$ a central charge
satisfying Bridgeland's support property (Bridgeland 2007
Definition~5.1). The local homeomorphism
$\mathrm{Stab}(X_5) \to \mathrm{Hom}(K_0^{\mathrm{num}}, \mathbb C)$
(Bridgeland 2007 Theorem~1.2) gives complex dimension $4$; mirror
symmetry lifts this to a cover of the stringy K\"ahler moduli. Neither
the local homeomorphism nor the $\widetilde{\mathrm{GL}}_2^+$-action
gives an $E_2$-monoidal category structure. The role of
$\mathrm{Stab}(X_5)$ is to parametrise the wall-crossing structure of
the Joyce--Brauer gerbe class of
Proposition~\ref{prop:cy-c-beyond-joyce-three-CY3}(i); the gerbe class
is itself stability-invariant (Bridgeland 2011 arXiv:1002.4374
Theorem~6.3).
\end{proof}

\begin{theorem}[$\Omega^{\mathrm{cat}}$ on the quintic recovers $X_5$ as non-commutative smooth proper]
\label{thm:cy-c-quintic-omega-cat-recovery}
\ClaimStatusProvedHere
The categorical cobar construction applied to
$B^{\mathrm{cat}}(\mathcal D(X_5))$ recovers $\mathcal D(X_5)$ up to
pro-equivalence:
\[
  \Omega^{\mathrm{cat}}\!\bigl( B^{\mathrm{cat}}(\mathcal D(X_5)) \bigr)
  \;\simeq\;
  \mathcal D(X_5)
  \qquad \text{in } \mathrm{Pro}(\mathrm{BrMon}^{\mathrm{comp}}_\infty).
\]
Via the Orlov-type reconstruction for non-commutative smooth proper dg
categories (Orlov 2016 arXiv:1402.7364, Theorem~3.28),
$\Omega^{\mathrm{cat}}$ recovers $X_5$ as a non-commutative smooth
proper variety with CY-$3$ structure: the dg category
$D^b_{\mathrm{dg}}(X_5)$ equipped with its $(-3)$-shifted symplectic
Hochschild cohomology.
\end{theorem}

\begin{proof}
Specialise Theorem~\ref{thm:cy-c-cat-thm-B} to
$\mathcal D = \mathcal D(X_5)$. The categorical Koszul locus condition
is verified by the CY-$3$ shifted symplectic structure on
$\HH^\bullet(\mathcal D(X_5))$ (Caldararu 2003): a smooth proper CY-$d$
dg category has a non-degenerate symmetric pairing of degree $(-d)$ on
Hochschild cohomology, giving the quadratic presentation required for
bar spectral sequence degeneration at $E_2$ (Keller--Van den Bergh 2011
arXiv:0804.3256, Theorem~4.4). The Pro-completion is necessary on the
strict stratum: the bar complex on a non-rigid target
(Proposition~\ref{prop:cy-c-beyond-quintic-fibre-functor}) is unbounded
above, and only the pro-object completion sees its contractibility
(the $(\infty,2)$-avatar of the Positselski co-derived enlargement,
Theorem~\ref{thm:cy-c-cat-thm-B}).

Orlov reconstruction: given a smooth proper dg category $\mathcal A$
with symmetric Hochschild pairing, the non-commutative variety
$X^{\mathrm{nc}}_{\mathcal A}$ whose derived category of coherent
sheaves recovers $\mathcal A$ is reconstructed up to equivalence by
the homotopy-orbit of the $\HH^0$-action (Orlov 2016 Theorem~3.28).
For $\mathcal A = D^b_{\mathrm{dg}}(X_5)$ the centre
$\HH^0 \simeq H^0(X_5, \cO_{X_5}) = k$ acts trivially, so the
reconstruction produces a non-commutative smooth proper CY-$3$
$X_5^{\mathrm{nc}}$ with
$D^b(\Coh(X_5^{\mathrm{nc}})) = D^b(\Coh(X_5))$. The upgrade to a
commutative variety is Ballard--Favero--Katzarkov 2014
(arXiv:1208.1682, Theorem~1.2): a smooth proper CY dg category with
commutative $\HH^0$ is equivalent to $D^b(\Coh(X))$ for a classical
commutative variety $X$ iff the compact generator reduces to a
coherent sheaf; on the quintic $\cO_{X_5}$ satisfies this, so
$X_5^{\mathrm{nc}} = X_5$.

The pro-inverse system realising $\eta_{\mathcal D(X_5)}$ is the
Postnikov-type tower of Orlov twists
$\mathrm{Tw}_n(D^b_{\mathrm{dg}}(X_5))$ indexed by $n \to \infty$,
converging pro-equivalently to $D^b_{\mathrm{dg}}(X_5)$ itself. The
tower is the $(\infty,2)$-analogue of Positselski's weight-completion:
at the algebra level the co-derived category sees the unbounded bar;
at the $(\infty,2)$-level the Pro-completion sees the unbounded
dg-tensor of Theorem~\ref{thm:cy-c-quintic-bcat-explicit}. Both
enlargements are required precisely because $X_5$ is strict CY$_3$
(no exceptional collection, no finite tilting object,
Proposition~\ref{prop:cy-c-quintic-kuznetsov-full}).
\end{proof}

\begin{corollary}[$\Phi_3^{\mathrm{cat}}(X_5)$ and the three categorical theorems]
\label{cor:cy-c-quintic-phi-three-cat}
\ClaimStatusProvedHere
The categorical CY-to-chiral functor $\Phi_3^{\mathrm{cat}}$ applied to
the quintic is
\[
  \Phi_3^{\mathrm{cat}}(X_5)
  \;:=\;
  B^{\mathrm{cat}}(\mathcal D(X_5))
  \;=\;
  \mathrm{Ind}\!\bigl(
    D^b_{\mathrm{dg}}(X_5) \otimes_{\mathrm{dgCat}_k}
    \mathrm{Fact}^{\mathrm{br}}_{X_5}
  \bigr)
  \;\in\;
  \mathrm{BrMon}^{\mathrm{comp}}_\infty.
\]
Under this identification:
\textup{(B)} $\Omega^{\mathrm{cat}}(\Phi_3^{\mathrm{cat}}(X_5))
  \simeq \mathcal D(X_5)$
in $\mathrm{Pro}(\mathrm{BrMon}^{\mathrm{comp}}_\infty)$
(Theorem~\ref{thm:cy-c-quintic-omega-cat-recovery});
\textup{(D)} $\kappa_{\mathrm{cat}}(\Phi_3^{\mathrm{cat}}(X_5))$
evaluates at the unit object to the scalar
$-25/3 = \chi_{\mathrm{top}}(X_5)/24 = -200/24$ (Euler number of the
quintic, Candelas--de la Ossa--Green--Parkes 1991), twisted by the
Joyce gerbe $[5 \cdot \mathrm{vol}] \in H^3(X_5, \mathbb C^\times)$;
\textup{(H)} $\mathrm{ChirHoch}^\bullet_{\mathrm{cat}}(\Phi_3^{\mathrm{cat}}(X_5))$
is supported on $\{0, 1, 2, 3\}$ with the degree-$3$ piece equal to
the Joyce gerbe (Corollary~\ref{cor:cy-c-cat-thm-H-tests}).
\end{corollary}

\begin{proof}
The explicit form of $\Phi_3^{\mathrm{cat}}(X_5)$ is
Theorem~\ref{thm:cy-c-quintic-bcat-explicit}. Readings~(B), (D), (H)
are Theorems~\ref{thm:cy-c-quintic-omega-cat-recovery},
\ref{thm:cy-c-cat-thm-D}, \ref{thm:cy-c-cat-thm-H} specialised to
$\mathcal D = \mathcal D(X_5)$ via
Definition~\ref{def:cy-c-quintic-dg-enhancement}. The scalar $-25/3$
is $\chi_{\mathrm{top}}(X_5)/24$; the Joyce gerbe value
$[5 \cdot \mathrm{vol}]$ is
Proposition~\ref{prop:cy-c-beyond-joyce-three-CY3}(i).
\end{proof}

\begin{remark}[Open thread: stability-refined $\Phi_3^{\mathrm{cat}, \sigma}$]
\label{rem:cy-c-quintic-open-stability}
Theorem~\ref{thm:cy-c-quintic-bcat-explicit} uses the factorisation
category on curves, independent of any Bridgeland stability choice.
A stability-refined $\Phi_3^{\mathrm{cat}, \sigma}(X_5)$ with
$\sigma \in \mathrm{Stab}(X_5)$ would incorporate central charges into
the braiding, yielding a family of braided monoidal $\infty$-categories
fibred over $\mathrm{Stab}(X_5)$. The wall-crossing automorphisms
should be the Joyce wall-crossing formulae lifted to the
$(\infty,2)$-level; the Joyce--Brauer gerbe class
$[\mathrm{ob}] \in H^3(\mathrm{Stab}, \mathbb C^\times)$ is the
obstruction to trivialising the family. Open: construct
$\Phi_3^{\mathrm{cat}, \sigma}$ and identify its wall-crossing with
the Kontsevich--Soibelman DT transformations (Kontsevich--Soibelman
2008 arXiv:0811.2435).
\end{remark}

\section{Explicit $B^{\mathrm{cat}}$ on the bicubic and the $(2,4)$-CI}
\label{sec:cy-c-beyond-bicubic-24CI-bcat-explicit}

The quintic construction of
Theorem~\ref{thm:cy-c-quintic-bcat-explicit} is driven by two
structural facts, neither of which uses the degree $5$: the ambient
is a smooth projective threefold with trivial canonical class, and
curves inside it have two-dimensional transverse normal bundle. Every
compact strict CY$_3$ complete intersection in $\mathbb P^{n+3}$
inherits both. This section states the construction for the two
remaining bicubic / $(2,4)$-CI members of the three-member hypersurface
family $(X_5, X_{3,3}, X_{2,4})$ and records the scalar values that
distinguish them.

\begin{proposition}[Kuznetsov component on the bicubic and the $(2,4)$-CI is the entire category]
\label{prop:cy-c-bicubic-24CI-kuznetsov-full}
\ClaimStatusProvedHere
For $X \in \{X_{3,3},\, X_{2,4}\}$ the Kuznetsov component
$\mathcal{K}u(X)$ coincides with $D^b(\Coh(X))$. Neither variety admits
a non-trivial semi-orthogonal decomposition, a tilting object of
finite global dimension, or an exceptional object.
\end{proposition}

\begin{proof}
Both $X_{3,3} \subset \mathbb P^5$ and $X_{2,4} \subset \mathbb P^5$ are
smooth projective CY$_3$ with $\omega_X \simeq \cO_X$ by the adjunction
formula applied to a bidegree-$(3,3)$ (resp.\ bidegree-$(2,4)$) complete
intersection: the ambient anticanonical $\cO_{\mathbb P^5}(6)$ cancels
against $\cO(3+3) = \cO(6)$ and $\cO(2+4) = \cO(6)$ respectively. The
Serre-duality argument of
Proposition~\ref{prop:cy-c-quintic-kuznetsov-full} applies verbatim:
an exceptional $E$ would satisfy
$\mathrm{Ext}^3(E,E) \simeq \mathrm{Ext}^0(E,E)^\vee = k$, contradicting
vanishing in positive degree. The Kuznetsov~2009
(arXiv:0904.4330, Lemma~2.5) trichotomy on admissible subcategories of
trivial-canonical smooth projective varieties collapses to the trivial
decomposition. The Bondal--Van den Bergh 2003
(arXiv:math/0204218, Theorem~3.1.1) classical generator is
$\cO_{X_{3,3}}$ (resp.\ $\cO_{X_{2,4}}$).
\end{proof}

\begin{proposition}[Joyce gerbe on the $(2,4)$-CI]
\label{prop:cy-c-beyond-joyce-24CI}
\ClaimStatusProvedHere
Let $X_{2,4} \subset \mathbb P^5$ be a smooth $(2,4)$ complete
intersection threefold: Hodge numbers
$(h^{1,1}, h^{2,1}) = (1, 89)$; $\chi_{\mathrm{top}}(X_{2,4}) = 2(1-89) = -176$;
$Y_3 = 2 \cdot 4 = 8$. The Joyce integration gerbe class is
\[
  \mathrm{ob}(X_{2,4}) \;=\; [8 \cdot \mathrm{vol}] \;\in\; J^3(X_{2,4})
  \;=\; \mathbb C^{90}/\Lambda,
\]
non-zero of infinite order.
\end{proposition}

\begin{proof}
Hodge numbers are Candelas--Font--Katz--Morrison 1994
(\emph{Nucl.\ Phys.\ B} 429, Table~1, arXiv:hep-th/9308083);
$\chi = 2(h^{1,1} - h^{2,1})$ for a simply connected CY$_3$ gives
$\chi = -176$. The Yukawa $Y_3 = H^3$ is computed on $\mathbb P^5$ by
Griffiths adjunction: the triple self-intersection of the hyperplane
class restricted to a complete intersection of bidegree $(2,4)$ is
$2 \cdot 4 = 8$. The Massey-product reduction of
Proposition~\ref{prop:cy-c-beyond-joyce-vanishing} applies as in case
(ii) of Proposition~\ref{prop:cy-c-beyond-joyce-three-CY3}: on a
strict compact CY$_3$ the triple Massey product on $H^1(T_X)$ reduces
to $Y_3$, and the image in $H^3(X_{2,4}, \mathbb C^\times)$ through
the intermediate Jacobian is non-zero of infinite order because
$\Omega^{3,0}$ is non-trivial. Joyce--Song 2012
(arXiv:0810.5645, Theorem~5.11) non-globalisation of the DT
integration map matches the Yukawa input.
\end{proof}

\begin{definition}[Factorisation category on curves in $X_{3,3}$, $X_{2,4}$]
\label{def:cy-c-bicubic-24CI-fact-cat}
\ClaimStatusDefinitional
For $X \in \{X_{3,3},\, X_{2,4}\}$, let $\mathrm{Curv}(X)$ denote the
stack of quasi-embedded smooth curves $\iota : C \hookrightarrow X$.
The braided factorisation $\infty$-category
\[
  \mathrm{Fact}^{\mathrm{br}}_X
  \;:=\;
  \lim_n \mathrm{QCoh}\bigl(\mathrm{Ran}_n(\mathrm{Curv}(X))\bigr)^{\otimes\text{-factorisable}}
\]
is defined by the Francis--Gaitsgory 2012 (arXiv:1110.5690, \S 3)
Ran-space construction. The $E_2$-structure arises from the
two-dimensional transverse normal bundle $N_{C/X}$ to a curve in a
threefold, by Ayala--Francis 2015 (arXiv:1206.5522, Theorem~1.2).
Genus-zero degree-$d$ components carry the Gromov--Witten / DT counts
$n_d(X_{3,3})$ and $n_d(X_{2,4})$ computed in Katz--Klemm--Vafa 1999
(arXiv:hep-th/9910181, \S 6) for bicubic and Klemm--Pandharipande
2008 (arXiv:math/0702189) for $(2,4)$-CI.
\end{definition}

\begin{theorem}[Explicit $B^{\mathrm{cat}}$ on the bicubic and the $(2,4)$-CI]
\label{thm:cy-c-bicubic-24CI-bcat-explicit}
\ClaimStatusProvedHere
For $X \in \{X_{3,3},\, X_{2,4}\}$ the categorical bar construction is
the braided presentable stable $\infty$-category
\[
  B^{\mathrm{cat}}(\mathcal D(X))
  \;=\;
  \mathrm{Ind}\!\bigl(
    D^b_{\mathrm{dg}}(X) \otimes_{\mathrm{dgCat}_k}
    \mathrm{Fact}^{\mathrm{br}}_X
  \bigr)
  \;\in\; \mathrm{BrMon}^{\mathrm{comp}}_\infty,
\]
obtained by the three-step procedure of
Theorem~\ref{thm:cy-c-quintic-bcat-explicit}: smooth-proper dg-tensor
(To\"en 2007 arXiv:math/0408337, Theorem~6.1), $E_2$-braiding from
the two-dimensional transverse Ran-space factorisation and the
$(-3)$-shifted symplectic Hochschild pairing (Caldararu 2003
arXiv:math/0308080, Theorem~5), $\mathrm{Ind}$-completion (Lurie
\emph{HA}~5.5.7.3). Compact generators are
$\cO_X \otimes \mathbb 1_C$ for $C \in \mathrm{Curv}(X)$.
\end{theorem}

\begin{proof}
The proof of Theorem~\ref{thm:cy-c-quintic-bcat-explicit} is
degree-independent: Step~(i) uses smooth-properness of
$D^b_{\mathrm{dg}}(X)$ (Orlov 2016 arXiv:1402.7364, Corollary~3.28)
and of the factorisation category on curves in a threefold
(Ben-Zvi--Francis--Nadler 2010 arXiv:0805.0157, Theorem~4.3.6), both
holding for any smooth projective CY$_3$. Step~(ii) uses only the
threefold assumption and the CY$_3$ Hochschild pairing, both present
on $X_{3,3}$ and $X_{2,4}$. Step~(iii) is formal.
Proposition~\ref{prop:cy-c-bicubic-24CI-kuznetsov-full} supplies the
CY$_3$ generator $\cO_X$ required in Step~(iii).
\end{proof}

\begin{theorem}[$\Omega^{\mathrm{cat}}$ on the bicubic and the $(2,4)$-CI recovers $X$]
\label{thm:cy-c-bicubic-24CI-omega-cat-recovery}
\ClaimStatusProvedHere
For $X \in \{X_{3,3},\, X_{2,4}\}$,
\[
  \Omega^{\mathrm{cat}}\!\bigl( B^{\mathrm{cat}}(\mathcal D(X)) \bigr)
  \;\simeq\; \mathcal D(X)
  \qquad \text{in } \mathrm{Pro}(\mathrm{BrMon}^{\mathrm{comp}}_\infty),
\]
and Orlov reconstruction (Orlov 2016 arXiv:1402.7364, Theorem~3.28)
recovers $X$ as a classical smooth projective CY$_3$ from the recovered
dg-enhancement $D^b_{\mathrm{dg}}(X)$.
\end{theorem}

\begin{proof}
Categorical Theorem~B
(Theorem~\ref{thm:cy-c-cat-thm-B}) specialises to
$\mathcal D = \mathcal D(X)$: the Koszul-locus hypothesis is the
CY$_3$ shifted-symplectic pairing on $\HH^\bullet(\mathcal D(X))$,
supplied by Caldararu 2003 (Theorem~5) for any smooth proper CY$_3$
dg category. Pro-completion is necessary because the bar is unbounded
on the strict stratum
(Proposition~\ref{prop:cy-c-bicubic-24CI-kuznetsov-full}). Orlov
reconstruction proceeds via $\HH^0(\mathcal A) = H^0(X, \cO_X) = k$
acting trivially; Ballard--Favero--Katzarkov 2014
(arXiv:1208.1682, Theorem~1.2) upgrades the non-commutative
reconstruction to the classical variety because $\cO_X$ is the
compact generator in both cases.
\end{proof}

\begin{corollary}[$\Phi_3^{\mathrm{cat}}$ on the bicubic and the $(2,4)$-CI]
\label{cor:cy-c-bicubic-24CI-phi-three-cat}
\ClaimStatusProvedHere
The categorical CY-to-chiral functor evaluates to
\[
  \Phi_3^{\mathrm{cat}}(X_{3,3})
  \;=\;
  \mathrm{Ind}\!\bigl(
    D^b_{\mathrm{dg}}(X_{3,3}) \otimes_{\mathrm{dgCat}_k}
    \mathrm{Fact}^{\mathrm{br}}_{X_{3,3}}
  \bigr),
\qquad
  \Phi_3^{\mathrm{cat}}(X_{2,4})
  \;=\;
  \mathrm{Ind}\!\bigl(
    D^b_{\mathrm{dg}}(X_{2,4}) \otimes_{\mathrm{dgCat}_k}
    \mathrm{Fact}^{\mathrm{br}}_{X_{2,4}}
  \bigr),
\]
both in $\mathrm{BrMon}^{\mathrm{comp}}_\infty$. The three categorical
theorems evaluate as follows:
\textup{(B)} $\Omega^{\mathrm{cat}}(\Phi_3^{\mathrm{cat}}(X)) \simeq \mathcal D(X)$
(Theorem~\ref{thm:cy-c-bicubic-24CI-omega-cat-recovery});
\textup{(D)} $\kappa_{\mathrm{cat}}(\Phi_3^{\mathrm{cat}}(X_{3,3})) = -144/24 = -6$
and $\kappa_{\mathrm{cat}}(\Phi_3^{\mathrm{cat}}(X_{2,4})) = -176/24 = -22/3$,
each twisted by the Joyce gerbe $[Y_3 \cdot \mathrm{vol}]$ with
$Y_3(X_{3,3}) = 9$ and $Y_3(X_{2,4}) = 8$;
\textup{(H)} $\mathrm{ChirHoch}^\bullet_{\mathrm{cat}}(\Phi_3^{\mathrm{cat}}(X))$
is supported on $\{0, 1, 2, 3\}$ with the degree-$3$ piece equal to the
Joyce gerbe class.
\end{corollary}

\begin{proof}
Explicit form from Theorem~\ref{thm:cy-c-bicubic-24CI-bcat-explicit}.
Reading~(B) from Theorem~\ref{thm:cy-c-bicubic-24CI-omega-cat-recovery}.
Reading~(D) from Theorem~\ref{thm:cy-c-cat-thm-D} specialised: the
scalar is $\chi_{\mathrm{top}}/24$, with the Joyce gerbe factor from
Proposition~\ref{prop:cy-c-beyond-joyce-three-CY3}(ii) ($X_{3,3}$,
$Y_3 = 9$) and Proposition~\ref{prop:cy-c-beyond-joyce-24CI}
($X_{2,4}$, $Y_3 = 8$). Reading~(H) from
Theorem~\ref{thm:cy-c-cat-thm-H} as in the quintic case.
\end{proof}

\begin{theorem}[Stratum-universal categorical scalar on compact CY$_3$ complete intersections in $\mathbb P^{n+3}$]
\label{thm:cy-c-ci-cy3-stratum-universal-kappa-cat}
\ClaimStatusProvedHere
Let $X \subset \mathbb P^{n+3}$ be a smooth complete intersection of
multidegree $(d_1, \ldots, d_n)$ with $\sum d_i = n+4$ (CY$_3$
condition) and $X$ on the strict stratum (no exceptional object, no
non-trivial SOD). The categorical scalar complementarity invariant is
\[
  \kappa_{\mathrm{cat}}(\Phi_3^{\mathrm{cat}}(X))
  \;=\; \chi_{\mathrm{top}}(X)/24,
\]
twisted by the Joyce integration gerbe
$\mathrm{ob}(X) = [Y_3(X) \cdot \mathrm{vol}] \in H^3(X, \mathbb C^\times)$
with Yukawa $Y_3(X) = d_1 d_2 \cdots d_n$. The three flagship members
of the hypersurface / CI stratum in $\mathbb P^{n+3}$ evaluate to the
table
\[
  \begin{array}{c|c|c|c|c}
    X & (h^{1,1}, h^{2,1}) & \chi_{\mathrm{top}} & \kappa_{\mathrm{cat}} & Y_3 \\ \hline
    X_5 \subset \mathbb P^4       & (1, 101) & -200 & -25/3   & 5  \\
    X_{3,3} \subset \mathbb P^5   & (1, 73)  & -144 & -6      & 9  \\
    X_{2,4} \subset \mathbb P^5   & (1, 89)  & -176 & -22/3   & 8
  \end{array}
\]
\end{theorem}

\begin{proof}
Theorem~\ref{thm:cy-c-cat-thm-D} gives the scalar reading of
$\kappa_{\mathrm{cat}}$ as the BCOV genus-one partition function
coefficient, which at one loop evaluates to $\chi_{\mathrm{top}}/24$
(Bershadsky--Cecotti--Ooguri--Vafa 1994, arXiv:hep-th/9309140,
eq.~(2.13); Costello--Li 2012, arXiv:1201.4501, Theorem~1.4). The
Yukawa formula $Y_3 = \prod d_i$ is the classical Griffiths adjunction
on a projective complete intersection. Hodge numbers for the three
entries: $X_5$ from Candelas--de la Ossa--Green--Parkes 1991 \S 4;
$X_{3,3}$ from Candelas--Font--Katz--Morrison 1994
(arXiv:hep-th/9308083, Table~1); $X_{2,4}$ likewise. Euler numbers
from $\chi = 2(h^{1,1} - h^{2,1})$ (simply connected, $h^{1,0} = h^{2,0} = 0$,
$h^{3,0} = 1$). None of the three admits exceptional objects by
Propositions~\ref{prop:cy-c-quintic-kuznetsov-full} and
\ref{prop:cy-c-bicubic-24CI-kuznetsov-full}; the strict-stratum
hypothesis is verified.
\end{proof}

\begin{remark}[Comparison: degree versus Euler along the hypersurface family]
\label{rem:cy-c-ci-cy3-degree-vs-euler}
The quintic / bicubic / $(2,4)$-CI triplet records a single mechanism
seen from three points of the complete-intersection moduli. The
Yukawa increases non-monotonically along degree-product
($5 \to 9 \to 8$), while the Euler number is also non-monotone
($-200 \to -144 \to -176$): the Yukawa counts holomorphic rigidity
through the triple intersection, the Euler number counts
Hodge-theoretic rigidity through $h^{2,1}$, and the two are
independent invariants of the CI multidegree. The categorical scalar
$\kappa_{\mathrm{cat}} = \chi/24$ and the gerbe scalar $Y_3$ are
distinct numerical shadows of the same stratum: the first is the
genus-one BCOV partition, the second is the classical intersection.
The quintic is the Yukawa-minimal entry ($Y_3 = 5$) and
Euler-most-negative ($\chi = -200$); the bicubic is Yukawa-maximal
($Y_3 = 9$) and Euler-least-negative ($\chi = -144$); the $(2,4)$-CI
sits between both. Bondal--Orlov reconstruction acts uniformly on
the three: no exceptional object, no non-trivial SOD, full category
is Kuznetsov component, classical variety recovered by Orlov's
$\HH^0 = k$ argument. The Joyce gerbe is non-zero of infinite order
in each case, with representative $[Y_3 \cdot \mathrm{vol}]$; the
categorical bar construction is uniform, differing only in the
scalar reading of Theorem~D.
\end{remark}

\begin{remark}[Open threads from the CI-CY$_3$ stratum]
\label{rem:cy-c-ci-cy3-open-threads}
Three threads remain open on the CI-CY$_3$ stratum.
\textup{(a)} Stability-refined $\Phi_3^{\mathrm{cat}, \sigma}$ for
$X_{3,3}$ and $X_{2,4}$: both carry finite-dimensional Bridgeland
stability manifolds
($\mathrm{rk}(K_0^{\mathrm{num}}(X)) = 4$ for a simply connected
CY$_3$ in projective space, Bayer--Macri 2011 arXiv:1103.5010), and
the wall-crossing is conjecturally controlled by the Joyce
integration gerbe, as for the quintic
(Remark~\ref{rem:cy-c-quintic-open-stability}).
\textup{(b)} CY$_3$ with $h^{1,1} > 1$: the general member of the
extended stratum admits a K\"unneth-like product SOD (e.g.\ bicubic in
$\mathbb P^2 \times \mathbb P^2$ with
$(h^{1,1}, h^{2,1}) = (2, 83)$, $\chi = -162$), which lies outside
the strict single-generator stratum of
Proposition~\ref{prop:cy-c-bicubic-24CI-kuznetsov-full} and requires
the Kuznetsov homological projective duality framework (Kuznetsov
2007 arXiv:math/0610957).
\textup{(c)} Non-compact local-mirror limits: large-volume limit
$t \to \infty$ in a one-parameter family degenerates to a
non-compact CY$_3$ with exceptional collection; the categorical bar
construction acquires finite fibre-functor shadow in that limit and
agrees with the quantum-group construction on the local-mirror side
(Aganagic--Okounkov 2017 arXiv:1704.08746, \S 2).
\end{remark}
